\documentclass[10pt,a4paper]{article}
% A standalone guide assembled from the explicit editorial section manifest.
% Compile with LuaLaTeX twice. Source IDs and archive links are retained.
\usepackage{fontspec}
\setmainfont{Latin Modern Roman}
\setsansfont{Latin Modern Sans}
\setmonofont{Latin Modern Mono}
\usepackage[english]{babel}
\babelprovide[import]{french}
\usepackage[a4paper,left=21mm,right=43mm,top=22mm,bottom=22mm,
  marginparwidth=29mm,marginparsep=5mm,headheight=14pt,headsep=5mm]{geometry}
\usepackage{amsmath,amssymb,mathtools}
\usepackage{newunicodechar}
\newunicodechar{∈}{\ensuremath{\in}}
\newunicodechar{⊗}{\ensuremath{\otimes}}
\newunicodechar{ø}{\ifmmode\text{\normalfont\itshape\char"00F8}\else\char"00F8\fi}
\newunicodechar{Ø}{\ifmmode\text{\normalfont\itshape\char"00D8}\else\char"00D8\fi}
\usepackage{xcolor}
\usepackage{tikz-cd}
\usetikzlibrary{shapes.geometric}
\usepackage{lua-ul}
\usepackage{marginnote}
\usepackage{fancyhdr}
\usepackage{pdflscape}
\usepackage{microtype}
\usepackage[unicode,breaklinks,colorlinks,linkcolor=blue!45!black,
  urlcolor=blue!45!black,pdfencoding=auto]{hyperref}
% Some transcribed symbols deliberately use the slashed-o letter. Select its
% actual text glyph even inside math; do not replace it by another math symbol.
% Install after package encoding declarations have initialized the text macros.
\AtBeginDocument{\RenewDocumentCommand{\o}{}{\ifmmode\text{\normalfont\itshape\char"00F8}\else\char"00F8\fi}}
\AtBeginDocument{\RenewDocumentCommand{\O}{}{\ifmmode\text{\normalfont\itshape\char"00D8}\else\char"00D8\fi}}
% Explicit maps also support minimal TeX Live installations.
\pdfextension mapfile{+cm.map}
\pdfextension mapfile{+symbols.map}
\pdfextension mapfile{+euler.map}
\pdfextension mapfile{+latxfont.map}
\pdfextension mapfile{+cmextra.map}
\pdfextension mapfile{+lm.map}
\setlength{\parindent}{0pt}
\setlength{\parskip}{4pt plus 1pt minus 1pt}
\setlength{\emergencystretch}{1.5em}
\setlength{\marginparpush}{3pt}
\raggedbottom
\definecolor{handblue}{RGB}{25,65,125}
\definecolor{deletedred}{RGB}{132,38,34}
\definecolor{strikegray}{RGB}{90,90,90}
\newcount\EditionMode
\EditionMode=2 % 0 diplomatic; 1 corrected apparatus; 2 reading
\tikzset{handwritten/.code={\ifnum\EditionMode=0
  \pgfkeysalso{/tikz/draw=handblue}\else\pgfkeysalso{/tikz/draw=black}\fi}}
\tikzcdset{handwritten/.style={/tikz/handwritten}}
\newif\ifInApparatus
\newif\ifEditionFrench
\EditionFrenchfalse
\newcommand{\EditionForm}{Guide to the manuscripts}
\newcommand{\CurrentScanURL}{}
\newcommand{\CurrentScanLabel}{}
\newcommand{\CurrentLinkPage}{}
\newcommand{\CurrentPageKey}{}
\pagestyle{fancy}
\fancyhf{}
\fancyhead[L]{\small\sffamily Nouveau EGA V}
\fancyhead[R]{\small\sffamily\CurrentScanLabel}
\fancyfoot[C]{\small\thepage}
\renewcommand{\headrulewidth}{0.2pt}
\newcommand{\ScanLink}[1]{\href{\CurrentScanURL}{#1}}
\newcount\SourceMarkSerial
\newcount\MarginNoteSerial
\NewDocumentCommand{\EditionMarginNote}{+m}{%
  \global\advance\MarginNoteSerial by 1
  \edef\MarginNoteIdentity{\CurrentPageKey:\the\MarginNoteSerial}%
  \leavevmode\hypertarget{marginanchor:\MarginNoteIdentity}{}%
  \marginpar{\raggedright\leavevmode\hypertarget{marginnote:\MarginNoteIdentity}{}#1}}
\newcommand{\SourceHref}[2]{%
  \global\advance\SourceMarkSerial by 1
  \hypertarget{egamark:#1:\CurrentLinkPage:\the\SourceMarkSerial}{\ScanLink{#2}}}
\newcommand{\scanpage}[1]{%
  \ifnum\EditionMode=2
    \def\NextScanPage{\ReadingScanPage{#1}}%
  \else
    \def\NextScanPage{\par\medskip\noindent{\small\sffamily\SourceHref{scan}{[scan #1]}}\par\smallskip}%
  \fi\NextScanPage}
\makeatletter
\newcommand{\ReadingScanPage}[1]{%
  \def\WaitingScanLabel{#1}%
  \@ifnextchar\mspage{\ReadingScanAndMsPage}{%
    \EditionMarginNote{\scriptsize\sffamily\SourceHref{scan}{scan #1}}}}
\def\ReadingScanAndMsPage\mspage#1{%
  \EditionMarginNote{\scriptsize\sffamily
    \SourceHref{scan}{scan \WaitingScanLabel}\par p.\,#1}}
\makeatother
\newcommand{\mspage}[1]{%
  \ifnum\EditionMode=2
    \EditionMarginNote{\scriptsize\sffamily p.\,#1}%
  \else
    {\small\sffamily[p.\,#1]}%
  \fi}
\NewDocumentCommand{\hand}{+m}{\ifnum\EditionMode=0{\color{handblue}#1}\else#1\fi}
\let\over\relax
\NewDocumentCommand{\over}{+m +m}{%
  {\color{strikegray}#1}\ensuremath{{}^{\nearrow}}{\color{handblue}#2}}
\NewDocumentCommand{\inserted}{O{} +m}{%
  {\color{handblue}\ensuremath{\langle}#2\ensuremath{\rangle}}}
\NewDocumentCommand{\deleted}{+m}{%
  {\color{deletedred}\ifmmode\text{[del.:}\,#1\,\text{]}\else[\textsf{del.:} #1]\fi}}
\NewDocumentCommand{\xout}{+m}{%
  {\color{strikegray}\ifmmode\text{[x:}\,#1\,\text{]}\else[\textsf{x:} #1]\fi}}
\NewDocumentCommand{\uline}{+m}{\ifmmode\underline{#1}\else\underLine{#1}\fi}
\NewDocumentCommand{\uuline}{+m}{\ifmmode\underline{\underline{#1}}\else\underLine[top=-1.05ex,bottom=-1.15ex]{\underLine{#1}}\fi}
\NewDocumentCommand{\unsure}{+m}{\SourceHref{unsure}{[#1?]}}
\newcommand{\illegible}[1]{\SourceHref{illegible}{\ifmmode\text{[illeg. #1]}\else[\textsf{illeg.\ #1}]\fi}}
\newsavebox{\MarginMeasure}
\newcount\MarginLineCount
\newcount\MarginContentSerial
% Count the actual packed rows of the margin-width paragraph. Do not count
% nested math boxes within a row as additional lines. An independent display
% or other box in its vertical list counts as a row of its own.
\directlua{
  function ega_margin_rows(boxnumber)
    local total = 0
    local hlist = node.id('hlist')
    local vlist = node.id('vlist')
    local function walk(n, vertical)
      while n do
        if n.id == hlist then
          if vertical then
            if n.width > 0 or n.height + n.depth > 0 then total = total + 1 end
          elseif n.list then walk(n.list, false) end
        elseif n.id == vlist and n.list then walk(n.list, true) end
        n = n.next
      end
    end
    walk(tex.box[boxnumber], false)
    return total
  end
}
\NewDocumentCommand{\MarginContent}{+m}{%
  \ifInApparatus
    {\small\textsf{\ifEditionFrench marg.\else marg.\fi:}\ #1}%
  \else
    \global\advance\MarginContentSerial by 1
    \ifmmode
      \typeout{EGA-MARGIN \CurrentPageKey\space\the\MarginContentSerial\space LINES -1 ROUTE inline-math}%
      \text{\scriptsize[\textsf{marg.:} #1]}%
    \else
      \sbox{\MarginMeasure}{\parbox[t]{\marginparwidth}{\footnotesize\raggedright #1}}%
      \directlua{tex.count['MarginLineCount'] = ega_margin_rows(\number\MarginMeasure)}%
      \ifnum\MarginLineCount<4
        \typeout{EGA-MARGIN \CurrentPageKey\space\the\MarginContentSerial\space LINES \the\MarginLineCount\space ROUTE margin}%
        \EditionMarginNote{\usebox{\MarginMeasure}}%
      \else
        \typeout{EGA-MARGIN \CurrentPageKey\space\the\MarginContentSerial\space LINES \the\MarginLineCount\space ROUTE body}%
        \par\smallskip\noindent\rule{16mm}{0.3pt}\par\nobreak
        \begingroup\small\noindent\textsf{\ifEditionFrench Note marginale\else Marginal note\fi.}\ #1\par\endgroup
        \nobreak\smallskip\noindent\rule{16mm}{0.3pt}\par\smallskip
      \fi
    \fi
  \fi}
\NewDocumentCommand{\margin}{+m}{{\color{handblue}\MarginContent{#1}}}
\NewDocumentCommand{\placedmargin}{+m}{\MarginContent{#1}}
\NewDocumentCommand{\diagram}{+m}{\par\smallskip{\centering #1\par}\smallskip}
\newenvironment{diagramblock}{\par\begin{center}}{\end{center}\par}
\newcommand{\VariantNumber}[1]{\csname variantnumber@#1\endcsname}
\newcommand{\VariantMark}[1]{%
  \ifInApparatus\else
    \ifmmode\text{\textsuperscript{\tiny\hyperlink{app:#1}{[\VariantNumber{#1}]}}}%
    \else\textsuperscript{\tiny\hyperlink{app:#1}{[\VariantNumber{#1}]}}\fi
  \fi}
\NewDocumentCommand{\var}{+m +m m m}{%
  \ifnum\EditionMode=0#1\else#2\ifnum\EditionMode=1\VariantMark{#4}\fi\fi}
\NewDocumentCommand{\tnote}{m +m}{%
  \ifnum\EditionMode=1\footnote{\textsf{Translator's note (#1).} #2}\fi}
\NewDocumentCommand{\restoredfrom}{m m +m}{%
  \begingroup\def\CurrentScanURL{#2}\def\CurrentLinkPage{#1}%
  \ifInApparatus
    {\small\sffamily\SourceHref{provenance}{[scan #1]}}\par
  \else
    \EditionMarginNote{\scriptsize\sffamily\SourceHref{provenance}{scan #1}}%
  \fi
  #3\par\endgroup}
\NewDocumentCommand{\AppEntry}{m m m +m +m +m}{%
  \par\smallskip\noindent\hypertarget{app:#1}{\textbf{[#2]}}\quad
  \textit{#3}. #4\par
  \begingroup\InApparatustrue\small
  \ifEditionFrench\textsf{Page :}\else\textsf{Page:}\fi\ 
  {\EditionMode=0 #5}\par
  \ifEditionFrench\textsf{Texte retenu :}\else\textsf{Retained text:}\fi\ 
  {\EditionMode=2 #6}\par\endgroup}
\newcommand{\AbsentReading}{\ifEditionFrench[omis]\else[omitted]\fi}
\newcommand{\WordSpaceReading}{\ifEditionFrench[espace entre les mots]\else[word space]\fi}
\NewDocumentCommand{\RegisterPage}{m +m +m}{%
  \expandafter\long\expandafter\gdef\csname PageBody@#1\endcsname{#2}%
  \expandafter\long\expandafter\gdef\csname PageApparatus@#1\endcsname{#3}}
\newcommand{\PrintPage}[3]{%
  \clearpage
  \global\MarginContentSerial=0
  \gdef\CurrentPageKey{#1}\gdef\CurrentScanLabel{#2}\gdef\CurrentLinkPage{#2}\gdef\CurrentScanURL{#3}%
  \typeout{EGA-SOURCE-BEGIN #1 PDF-PAGE \thepage}%
  \pdfbookmark[1]{Scan #2}{source:#1}%
  \csname PageBody@#1\endcsname\par
  \ifnum\EditionMode=1
    \csname PageApparatus@#1\endcsname
  \fi
  \typeout{EGA-SOURCE-END #1 PDF-PAGE \thepage}}

\begin{document}
\selectlanguage{english}
\hypersetup{pageanchor=false}
\begin{titlepage}\thispagestyle{empty}\vspace*{20mm}
{\Large Alexandre Grothendieck\par}\vspace{12mm}
{\Huge\bfseries Nouveau EGA V\par}\vspace{8mm}
{\Large Guide to the manuscripts\par}\vspace{7mm}
{\large French prenotes and Piotr Blass's English typescript\par}\vspace{12mm}
Université de Montpellier, fonds Grothendieck\\ Cotes 25 and 24\\
14 November 2026
\vfill\end{titlepage}\clearpage\hypersetup{pageanchor=true}
\section*{Provenance, credit, and rights}
AI-transcribed by Claude Opus 5.5 (Anthropic, effort high, Message Batches API); checked against the scans, and this edition built, by GPT-6 Astra (OpenAI, ChatGPT Pro mode); prompts and records by Claude (Anthropic).

Guide written by GPT-6 Astra (OpenAI, ChatGPT Pro mode).

Cote 25 is described in the archive as dated “[à partir de 1967-1987]”. Cote 24 is Piotr Blass's translation typescript dated 1987.

The translation in cote 24 is Piotr Blass's (1987); the rights in it belong to its translator or his heirs.

The Grothendieck Circle scan EGA5.pdf, the P. and J. Blass translations (2004), J. C. Reinhold's web re-edition, and the Centre for Grothendieckian Studies item L49d are published texts not reproduced and not consulted for any reading. Published EGA references used for mathematical context are identified in the guide's editorial discussions and reference section.

The rights in the original work belong to its author or the author's representatives or heirs. Where we hold rights, in the AI-produced transcription, translation, corrections and typesetting, they are dedicated to the public domain under CC0 1.0. This edition is published for educational, non-profit purposes, and we claim no rights we do not hold. If you hold rights in this work and object to its publication here, please open a GitHub issue and we will take it down.

\clearpage\gdef\CurrentScanLabel{Guide}
\par\typeout{EGA-SOURCE-BEGIN guide-iv16 PDF-PAGE \thepage}
\pdfbookmark[1]{IV 16: singular and supersingular loci}{source:guide-iv16}
\paragraph{IV 16, definitions 1--4 and the intended placement (cote 25, sites 139--141).}
The packet is headed ``Ensemble singulier et supersingulier d'une fonction, et critères différentiels'' and expressly says that the new number will be used in the paragraph on hyperplane sections, while its natural place is paragraph 16. Site 139 is an exceptionally faint carbon impression: its numerous doubtful readings remain independent of the clearer page that follows it. On site 140, definition 1 starts from a regular prescheme, a function $\phi$, and the class of $\phi_x$ in $\mathfrak m_x^2/\mathfrak m_x^3$. It defines the singular and supersingular zeros by the order-two condition and the degeneracy of the resulting quadratic form; the word ``non'' before ``dégénérée'' is deleted. The same page calls a singular zero that is not supersingular an ordinary singular zero. Remark 2 discusses the relation with regularity of $V(\phi)$, then a substantial passage involving multiplicity is crossed out. Definition 3 introduces the geometric version over a field. Its actual typing has $k'$ in ``extension $k'$ de $k$'', but the following line has $k$ followed by a low handwritten comma, not a raised prime; the reading preserves both. Remarks 4 on site 141 discuss the field-valued test, the relative version fiber by fiber, and invariance under multiplication of $\phi$ by a unit, which motivates extending the definitions to a section of an invertible Module. Several base letters there remain overwritten or uncertain between $S$ and $Y$. These definitions, corrections and uncertainties can be checked at \href{https://grothendieck.umontpellier.fr/web/viewer.html?file=../25.pdf#page=140}{site 139}, \href{https://grothendieck.umontpellier.fr/web/viewer.html?file=../25.pdf#page=141}{site 140} and \href{https://grothendieck.umontpellier.fr/web/viewer.html?file=../25.pdf#page=142}{site 141}.

\paragraph{IV 16, principal parts, discriminant and propositions 5--7 (sites 141--148).}
The construction proceeds through the sections $d^i\phi$ of the sheaves of principal parts. The published reference for principal parts is \href{https://www.numdam.org/articles/10.1007/BF02732123/}{EGA IV 16.3.1--16.3.2}; the use of the diagonal there gives the context for these differential constructions. Proposition 5 identifies the zero and singular-zero loci using orders zero and one; its proposed order-two clause is deleted. The surrounding rectangle on site 141 carries the instruction to keep the adjacent text, so the rectangle itself does not delete its contents. Site 142 defines the closed singular subscheme from the first differential and obtains a quadratic section $M(\phi)$ on it from the second principal part. The manuscript introduces the corresponding map from the tangent Module to the differential Module and the determinant $D(\phi)$; site 143 then designates the closed subscheme cut out by this determinant as the supersingular subscheme and leaves a bare heading for proposition 6, asking that the construction be collected there. The next argument assumes the stated smoothness and relative-dimension hypotheses on $X/Y/S$ and on $V^1=V(\phi)^{\mathrm{sing}}$. Its diagram (D) is built from the transitivity sequence and the conormal/differential sequence; site 144 explicitly refers the latter to paragraph 16 and invokes the Jacobian criterion. Proposition 7 on site 145 states the description of the supersingular locus by ramification of $V^1\longrightarrow Y$. The following argument seeks the corresponding identity of closed subschemes. It compares the ideals attached to two maps in a diagram with local direct factors $P,Q\subset M$, ranks satisfying $p+q=m+1$, and the determinant line $L=\det P\otimes\det Q\otimes\det M^{-1}$. Site 146 supplies the square with its single $L$ object and the two maps $a,b$; the text explicitly leaves a possible sign dependent on conventions. Site 147 identifies the relevant component of $a$ with the earlier discriminant, again with that sign reservation, and asks for a recapitulative proposition for a section of an invertible Module. Site 148 supplies the wholly handwritten diagram (D') for that extension. The author also discusses isolation of ordinary singularities, refers to the Cohen--Macaulay property and excellence, and asks where the isolation result should be put in Chapter IV. The argument's written coefficients and indices are retained: site 144 has $\mathcal O_X$ in the first displayed source, says that $W$ is smooth over $X$, and repeats $\Omega^1_{X/S}$ at the end of the subsequent conormal display; its principal-parts decomposition and quadratic expression use $X/S$ where diagram (D) uses $X/Y$. Site 145 changes rank letters between the exterior-power expressions and has a visible $M$ in the displayed Hom target. Site 146 leaves the connecting signs blank in the sample exact sequence and determinant identity. Site 147 lacks the dual sign on the later intersection's $P^1$ and writes $\Omega^1_{X/Y}+\Omega^1_{Y/S}$ in its second exact sequence. Site 148 prints the dimension thresholds in terms of $\dim_x X$ and immediately gives the plane example $\phi=xy$; the edition preserves this juxtaposition. The tensor exponent on $K$ in diagram (D') remains doubtful, and the horizontal endpoints are solid dots while the vertical endpoints are open zeros. These are observations of the written argument, not completed corrections. See \href{https://grothendieck.umontpellier.fr/web/viewer.html?file=../25.pdf#page=143}{sites 142--143}, \href{https://grothendieck.umontpellier.fr/web/viewer.html?file=../25.pdf#page=145}{sites 144--146} and \href{https://grothendieck.umontpellier.fr/web/viewer.html?file=../25.pdf#page=148}{sites 147--148}.

\paragraph{IV 16, augmented rings and the passage to several equations (sites 149--152).}
The author now explicitly says that one equation is still too particular and proposes an homomorphism $E\longrightarrow\mathcal O_X$ with $E$ locally free of rank $m$. Site 149 introduces a regular augmentation $A\longrightarrow A/I$ by projectivity of $I/I^2$ and the symmetric-algebra comparison with $\mathrm{gr}_I(A)$, while asking whether the terminology adopted in paragraph 19 requires ``quasi-régulière''. The terminology can be compared with \href{https://www.numdam.org/articles/10.1007/BF02732123/}{published EGA IV 16.9.1--16.9.3}, which defines regular and quasi-regular ideals and immersions and uses the symmetric-algebra comparison in 16.9.3. It then defines an augmentation of ordinary quadratic type through a nondegenerate invertible relation Module in degree two. Proposition 8 and lemma 8.1 on site 150 formulate the graded projectivity and rank criteria and discuss multiplication by the quadratic element; the author observes that the lemma's proof uses nonvanishing in the reduced fibers and suggests a more general formulation. The edition preserves both occurrences of the displayed binomial difference with lower index $d$, while the rank parameter is also called $d$; it also preserves the later wording that calls $B$ ``le spectre d'un corps'' after introducing it as a ring. Proposition 9, beginning on site 150 and continuing on site 151, treats the quotient by the image of a projective Module and its induced augmentation. It expresses the condition through the kernel of the first-order map and the induced quadratic map, then writes out how this latter map is obtained from multiplication in the associated graded algebra. One formula runs beyond the right image edge and is not completed editorially. Corollary 10 passes to a regular local ring and gives the formulations in terms of successive equations, the first-order rank and a nondegenerate quadratic element. Its continuation on site 152 explicitly refers to definition 1. The text then introduces singular, ordinary singular and supersingular zeros for the homomorphism $u$, using $\mathcal O_X$-regularity and the local ring of $V(u)$, and mentions the geometric versions over a field. The underlined capital $M$ at the target of $u$ in corollary 10 is retained as written. The long diagonal strokes across this packet have an unresolved cancellation or working-mark function: site 150's stroke stops before proposition 9, site 151 has two separate strokes, and site 152's stroke stops above the later terminology discussion. The text remains present with visible uncertain marginal marks, and the local cancellations remain separately encoded. See \href{https://grothendieck.umontpellier.fr/web/viewer.html?file=../25.pdf#page=150}{site 149}, \href{https://grothendieck.umontpellier.fr/web/viewer.html?file=../25.pdf#page=151}{site 150}, \href{https://grothendieck.umontpellier.fr/web/viewer.html?file=../25.pdf#page=152}{site 151} and \href{https://grothendieck.umontpellier.fr/web/viewer.html?file=../25.pdf#page=153}{site 152}.

\paragraph{IV 16, the differential construction for a Module and its unfinished continuation (sites 153--156).}
Site 153 explains the dependence on the local ring of $V(u)$ apart from the dimension/regularity condition, discusses the rational-point and relative versions, and notes the equivalent presentation of a section through the dual Module. It constructs $V^1\subset V^o=V(u)$ from the maximal exterior power of the first differential. Site 154 explicitly applies proposition 9, reduces the ordinary-singularity test to order two, and introduces the further rank-drop locus $V''$, the kernel line $K$ and cokernel $W$ on $V^1-V''$, the quadratic map $K\longrightarrow\mathrm{Sym}^2(W)$, and its discriminant $D$. Site 155 puts $V^2=V(D)$, records the ordinary-singular locus as $V^1-V''-V^2$, and interposes a local-freeness lemma to justify the construction. Its other displayed expression, for the supersingular locus, contains an actual blank between $V''$ and $V^2$; no union or difference sign has been supplied. The lemma begins with the letter $v$ and subsequently uses $u$. The surrounding pages repeatedly replace typed $S$ by handwritten $Y$, while other occurrences of $S$ remain untouched; site 154 changes rank letters from $n$ to $d$ but leaves $n$ in the discriminant exponent, whose later sign remains overwritten and uncertain. These incomplete revisions are preserved. The final discussion asks for generalizations of propositions 6 and 7, stipulates the smoothness and dimension data for $V^o$ and $V^1$, and explains the intended minimum-dimension choice. On site 156 the proposed comparison with the ramification subscheme is expressly left to be verified: the author says he has not carried out the exercise in detail and points back to the dévissage after proposition 7. The final paragraph records a rank-locus codimension assertion, while its proposed induction explanation is canceled. The guide supplies neither that missing verification nor a new theorem. See \href{https://grothendieck.umontpellier.fr/web/viewer.html?file=../25.pdf#page=154}{site 153}, \href{https://grothendieck.umontpellier.fr/web/viewer.html?file=../25.pdf#page=155}{site 154}, \href{https://grothendieck.umontpellier.fr/web/viewer.html?file=../25.pdf#page=156}{site 155} and \href{https://grothendieck.umontpellier.fr/web/viewer.html?file=../25.pdf#page=157}{site 156}.
\par\typeout{EGA-SOURCE-END guide-iv16 PDF-PAGE \thepage}

\par\typeout{EGA-SOURCE-BEGIN guide-iv17-15 PDF-PAGE \thepage}
\pdfbookmark[1]{IV 17.15: smooth forms}{source:guide-iv17-15}
\paragraph{IV 17.15, the terminological revision of SUP (cote 25, sites 157--160).}
Site 157 carries an exceptionally faint forward-reading impression; it is retained with extensive local doubts rather than treated as blank. The clearer opening on site 158 is headed ``Formes lisses'', with ``augmentations et singularités élémentaires'' inserted and the earlier continuation ``formes quadratiques lisses'' deleted. Its opening note explicitly revisits the preceding manuscript, now called SUP. It recalls the correction of ``non dégénérée'' to ``dégénérée'' in definition 1 and discusses the characteristic-two example $x^2+yz=0$. The author uses that example to explain his proposed vocabulary for smooth quadratic forms, nondegenerate quadratic forms, ordinary or elementary quadratic singularities, and supersingularity. The text states its relation between smoothness and nondegeneracy, including the rank and residual-characteristic qualifications, and asks the recipient's opinion about replacing ``ordinaire'' by ``élémentaire'' and extending the term to arbitrary multiplicity. Site 159 says that proposition 7 of SUP is intended to retain supersingularity in its statement, proposes a further closed subscheme of ultra-singular zeros using the ``discriminant corrigé'', and expressly requires the terminology on SUP page 10 to be revised as well. It also identifies IV 23.9.2 and the conditions of 8.8 as needing a change. The ensuing claim about the inseparability degree on site 160 is qualified by the author's own ``je crois'' and ``sans garantie''; the guide retains that status and supplies no missing verification. The overwritten letters, local cancellations and inserted terms remain visible in the transcription. See \href{https://grothendieck.umontpellier.fr/web/viewer.html?file=../25.pdf#page=158}{site 157}, \href{https://grothendieck.umontpellier.fr/web/viewer.html?file=../25.pdf#page=159}{site 158}, \href{https://grothendieck.umontpellier.fr/web/viewer.html?file=../25.pdf#page=160}{site 159} and \href{https://grothendieck.umontpellier.fr/web/viewer.html?file=../25.pdf#page=161}{site 160}.

\paragraph{IV 17.15.1--15.2, the definitions and their written dependencies (sites 160--162).}
Number 15.1 starts from a finite locally free Module $E$ on $S$ and a section $\phi$ of $\mathrm{Sym}^n(E)$, identifies it with a section of $\mathcal O_{P(E)}(n)$ by an explicit reference to Chapter III, and defines a smooth form through smoothness of $V(\phi)$ together with the stated nonvanishing condition on every fiber. It discusses fiberwise testing and change of the ground field, then gives the polynomial criterion in terms of $\phi$ and its partial derivatives, expressly invoking the Jacobian criterion. The inequality on $n$ is overwritten; the small inserted condition near the criterion is still doubtful, and the margin separately observes that degree zero gives the invertible sections of $\mathcal O_S$. The written rank exponent of $E$ and the variable list $X_0,\ldots,X_r$ use different letters and have been retained. Number 15.2 defines an elementary augmentation of multiplicity $n$ by three stated requirements: local freeness of $J/J^2$, generation of the graded kernel by an invertible Module locally a direct factor in degree $n$, and smoothness of its local generating form. Its displayed map is the symmetric-algebra comparison with $\mathrm{gr}_J(\mathcal O_X)$; $K$ is added as the name of the degree-$n$ kernel, and $\mu$ is handwritten over an obscured earlier letter. The accompanying note proposes formulating the last two requirements through a smooth invertible submodule, referring back to 15.1. The text then defines the corresponding language for augmented rings, local rings, points of a prescheme and geometric singularities over a field, and asks for a remark connecting the multiplicity with Samuel's notion. Numbers 15.3 and 15.4 request the multiplicity-$n$ variants of SUP proposition 8, proposition 9 and corollary 10; they provide instructions for further writing rather than new proofs here. Number 15.5 similarly requests the definition for a section of the structure sheaf or a locally free Module. See \href{https://grothendieck.umontpellier.fr/web/viewer.html?file=../25.pdf#page=161}{site 160}, \href{https://grothendieck.umontpellier.fr/web/viewer.html?file=../25.pdf#page=162}{site 161} and \href{https://grothendieck.umontpellier.fr/web/viewer.html?file=../25.pdf#page=163}{site 162}.

\paragraph{IV 17.15.6--15.7, local and geometric descriptions (sites 162--164).}
Number 15.6 writes a local description inside a regular locally noetherian ambient prescheme, using the stated regular sequence and the elementary singular zero of its defining functions. It then records the complete-intersection consequence and appeals to 15.4 for the reduction to one equation in a regular intermediate prescheme. The continuation on site 163 states the Cohen--Macaulay consequence and postpones the isolation assertion to a blow-up calculation in paragraph 20. The normality and reducedness thresholds and the zero-dimensional description are preserved as written. The zero-dimensional discussion has a handwritten replacement of ``dim'' by ``rang'', specifies the least exponent annihilating the maximal ideal, and continues on site 164 with the polynomial quotient over the ground field. Number 15.7 gives the corresponding geometric formulation, referring expressly to 15.6 and 15.5; it describes the value of the $n$th principal part in the degree-$n$ augmentation piece and the smoothness of the resulting form. Its residue-field terms bear overlapping corrections that remain doubtful. The author inserts a bracketed reminder that openness of the smoothness locus of a form was omitted from 15.1. The end of the argument discusses geometric normality and reducedness, the rational zero-dimensional case, and the classical description of a plane singularity with distinct tangents. The following parenthesis asks that the agreement of the absolute and relative notions at a rational point be stated already in 15.2 and discusses a separable residue extension, with the author's own reservation about the word ``finie''. The old canceled version of the opening of 15.7 survives between the new typed lines on site 163 and remains separately transcribed. See \href{https://grothendieck.umontpellier.fr/web/viewer.html?file=../25.pdf#page=163}{site 162}, \href{https://grothendieck.umontpellier.fr/web/viewer.html?file=../25.pdf#page=164}{site 163} and \href{https://grothendieck.umontpellier.fr/web/viewer.html?file=../25.pdf#page=165}{site 164}.

\paragraph{IV 17.15, the proposed rearrangement at the close (site 164).}
The final NB recalls the short paragraph 20 discussion of blowing up along a closed subscheme whose augmentation is elementary of multiplicity $n$. It identifies its ingredients as the general smoothness results of paragraph 17 and the beginning of 17.15.2, then adds that even the latter definition might be avoided. It proposes moving these results immediately after 15.1 and grouping the geometric part of 15.2 and the subsequent material into 17.16. The interlinear phrase concerning that geometric material is typed and placed before the surviving range 15.3--15.7; the range itself is not canceled, and its overwritten digit remains uncertain. The paragraph concludes that the planned number 15 has expanded into at least two parts: general statements about elementary objects of arbitrary multiplicity, and the characterization of smooth quadratic forms. It expressly says that the latter borrows only 15.1, practically the definition of a smooth form. This is the manuscript's proposed organization, with its changing numbering preserved. The additions ``du début de'' and ``au moins'' are typed and positioned by handwritten leaders; the last sentence has a pen correction from ``contient'' to ``contenant''. See \href{https://grothendieck.umontpellier.fr/web/viewer.html?file=../25.pdf#page=165}{site 164}.
\par\typeout{EGA-SOURCE-END guide-iv17-15 PDF-PAGE \thepage}

\par\typeout{EGA-SOURCE-BEGIN guide-sup-corrections PDF-PAGE \thepage}
\pdfbookmark[1]{Corrections and invariance of SUP}{source:guide-sup-corrections}
\paragraph{Corrections to ``Ensemble singulier et supersingulier'' (cote 25, sites 165--167).}
The faint opening on site 165 and the clearer page headed ``EGA IV par ?'' on site 166 introduce a further set of corrections to the earlier notes. Site 165 retains extensive local doubts; its pale impression is not silently replaced by the clearer wording on site 166. The opening instructions on site 166 change $S$ to $Y$ on manuscript pages 14 and 15 and object to the ``i.e.'' on page 15, line 10. The subsequent argument addresses the generalization of proposition 7 mentioned on pages 16 and 17. Its diagram combines the conormal sheaf of $V^1$ in $X$ with the relative differential sheaves for $X/Y$, $X/S$, $Y/S$ and $V^1/S$. It explicitly invokes the lemma on manuscript page 7 and compares the zero scheme $V(a)$ with $V(D)$. The nested immersions $V^1\subset V^o\subset X$ provide the displayed conormal sequence, for which the text leaves a reference to be supplied. The two following identifications use $E$ and $\mathrm{Hom}(W,K)$; $W$ and $K$ are described by the cokernel and kernel of the displayed map $v^1$. Site 167 refers back to manuscript page 5 for the first identification and to the attached differential complements to ``variation du rang d'un homomorphisme de Modules loc libres'' for the second. It then separates two checks: identifying $\mu i$ with $v^1$, and identifying the induced $\mu'$ with the homomorphism of the quadratic form from the bottom of page 15. The author leaves these checks to Dieudonné and expressly raises a possible sign. The guide therefore records the route of the proposed verification and its dependencies, without supplying that sign or claiming that the omitted checks were performed. The written comparison signs, the occurrences of $V(u)$, and the unsettled overwritten cells remain in the text. See \href{https://grothendieck.umontpellier.fr/web/viewer.html?file=../25.pdf#page=166}{site 165}, \href{https://grothendieck.umontpellier.fr/web/viewer.html?file=../25.pdf#page=167}{site 166} and \href{https://grothendieck.umontpellier.fr/web/viewer.html?file=../25.pdf#page=168}{site 167}.

\paragraph{The proposed intrinsic descriptions and their unfinished verification (sites 168--171).}
Site 168 revisits the change of viewpoint mentioned in the author's last letter. It starts with a closed subscheme $V^o=V(J)$ of $X$, the conormal Module $N=J/J^2$, and a local presentation through a finite locally free Module $E$. The discussion asks for the independence of $V^1$, $V''$ and $V^2$ from that presentation, while retaining the dependence on the written ambient data and on $m$. It describes $V^1$ and $V''$ by the exterior-power maps labelled $(*_m)$ and $(*_{m-1})$. The opening formula on site 169 has an unmatched parenthesis in the scan, which has been retained. For $V^2$, site 169 expressly says that the invariance has not really been verified and gives it conjectural status. The proposed route compares the maps from $E$ and from $N$, names their common cokernel $W$, and introduces the kernel $k$ of the intrinsic map as a quotient of $K$. It asks for the previously defined $M:K\to\mathrm{Sym}^2(W)$ to factor through a presentation-independent $m:k\to\mathrm{Sym}^2(W)$. The rank phrase beside $W$ contains a disputed cancellation and is left unresolved. Site 170 continues this proposed route with a representing functor, whose slashed-loop symbol remains locally ambiguous, then describes the intended complementary locus in terms of flatness, geometric regularity or ordinary quadratic singularities, and transverse regularity. The displayed prose also has an actual blank between $E$ and $J$; an arrow has not been supplied. Site 171 discusses the inverse-image situation, the available global surjection from $E^o$ to the conormal Module, and the resulting discriminant expression. Its handwritten insertion before ``à priori'' reads ``ici''. The final paragraph proposes adding the invariance discussion in small print after the locally free Module case. It weighs paragraph 19 against paragraph 22, explicitly locating the differential argument used on manuscript pages 18 and 19 with the Grassmannian material, and suggests a matrix calculation as an alternative exposition. These are the author's stated proposals and reservations. See \href{https://grothendieck.umontpellier.fr/web/viewer.html?file=../25.pdf#page=169}{site 168}, \href{https://grothendieck.umontpellier.fr/web/viewer.html?file=../25.pdf#page=170}{site 169}, \href{https://grothendieck.umontpellier.fr/web/viewer.html?file=../25.pdf#page=171}{site 170} and \href{https://grothendieck.umontpellier.fr/web/viewer.html?file=../25.pdf#page=172}{site 171}.
\par\typeout{EGA-SOURCE-END guide-sup-corrections PDF-PAGE \thepage}

\par\typeout{EGA-SOURCE-BEGIN guide-remords PDF-PAGE \thepage}
\pdfbookmark[1]{Remords on transversality}{source:guide-remords}
\paragraph{The ``Remords'' on transversality (cote 25, site 172).}
This one-page note is headed ``EGA IV 19'' and ``Remords à Ensemble singulier et supersingulier''. It explicitly follows the author's drafting of the transversality sorite and reconsiders the setting of the earlier singularity notes. The opening starts from a section $u$ of $W(\underline E)$ and identifies its zero scheme with the inverse image of the zero section. The interlinear phrases ``sous-préschéma $E_o$ de $E$'' and ``pour u'' are typed additions positioned by pen leaders. The text then proposes treating a general $X$-prescheme $E$, a subscheme $E_o$ and a section, with inverse image $X_o$. Its added smoothness sentence reads ``en supposant'', with pen retouches inside otherwise typed text. The author expressly reserves for detailed writing the question of an additional relative-dimension restriction on $E_o$. The proposed terminology for an ordinary point of nontransversality is written above the canceled word ``singulier''; the page relates it to a geometric ordinary quadratic point of $X_{oy}$ and the displayed codimension requirement. The following conormal expression uses $v_o^*(N_{E_o/E})$ in the role earlier occupied by $E\otimes O_{V(u)}$. The transitions between $u$ and $v$, the earlier ``rel. à S'', the literal spellings and the local overprints have all been retained. At one replacement there are two uncancelled $x$ marks, one typed above-left and one handwritten above-right of the canceled expression; the scan does not establish which should supersede the other. The handwritten abbreviation before ``transversalement'' reads ``un imm.'' and is not expanded. The ending asks the recipient to write in this setting, to give the vector-bundle translations at the end, and to discuss any initial difficulties during the next Paris visit. The added final sentence promises ``des notes succintes ces jours-ci''. This guide describes the proposed revision and its stated reservations; it does not supply the detailed verification requested by the author. See \href{https://grothendieck.umontpellier.fr/web/viewer.html?file=../25.pdf#page=173}{site 172}.
\par\typeout{EGA-SOURCE-END guide-remords PDF-PAGE \thepage}

\par\typeout{EGA-SOURCE-BEGIN guide-iv17-16 PDF-PAGE \thepage}
\pdfbookmark[1]{IV 17.16: quadratic forms}{source:guide-iv17-16}
\paragraph{IV 17.16: smooth quadratic forms (cote 25, sites 173--184).}
Site 173 contains an extremely faint forward-reading impression of the opening; its recovered text remains qualified throughout. Site 174 carries the clearer heading ``17.16. Formes quadratiques lisses'' and begins with 16.1. A section \(Q\) of \(\underline{\mathrm{Sym}}^2(E)\) gives a symmetric bilinear form on the dual Module, a homomorphism \(\check E\longrightarrow E\), and then a map of determinant Modules. The resulting section \(d(Q)\) is initially called the discriminant, with the author's express reservation about that name. Proposition 16.2 gives the first smoothness criterion, organized by the residue characteristic and parity of the rank. Its proof uses partial derivatives, the displayed identity \(2Q=\sum_i X_i\,\partial Q/\partial X_i\), and the kernel of the alternating bilinear form in characteristic two. On site 175 the author criticizes this use of coordinates and isolates Lemma 16.3 as the intended elementary calculation concerning a point of \(V(Q)\) and the kernel of \(\check E\longrightarrow E\). The displayed \(Q(x,x)\), the printed rank expression, and the altered letters in the derivative formulas remain as written. Paragraph 16.4 then identifies the odd-rank, characteristic-two case as the part motivating the following study and writes the even- and odd-rank standard forms explicitly. See \href{https://grothendieck.umontpellier.fr/web/viewer.html?file=../25.pdf#page=174}{site 173}, \href{https://grothendieck.umontpellier.fr/web/viewer.html?file=../25.pdf#page=175}{site 174}, and \href{https://grothendieck.umontpellier.fr/web/viewer.html?file=../25.pdf#page=176}{site 175}.

Lemma 16.5 begins on site 175 with the universal quadratic polynomial in the coefficients \(a_i,b_{ij}\) and continues on site 176 with the content of its determinant polynomial. The manuscript then introduces \(\widetilde d(Q)=\tfrac12d(Q)\) in the odd-rank discussion. Proposition 16.6 formulates the resulting section over a prescheme; its condition on base change receives a handwritten addition about functoriality under isomorphisms. Site 177 follows this with an explicit editorial objection to stating base-change compatibility before that functoriality, since the base-changed Module is specified up to isomorphism. Definition 16.7 names the corrected discriminant, and Corollary 16.8 gives the two written direct-sum formulas, including the handwritten factor 2 in its second case. The relation in 16.6(c) has a visibly overworked raised 2 after \(\widetilde d(Q)\); the transcription retains that uncertainty alongside the preceding half-determinant formula. Theorem 16.9 lists smoothness, invertibility of the corrected discriminant, and the standard-form descriptions after the indicated surjective and faithfully flat base changes. Corollary 16.10 on site 178 gives the field version and discusses the algebraically closed case. These numbered statements are presented with their actual revisions and disputed symbols, rather than with a silently reconciled formula. See \href{https://grothendieck.umontpellier.fr/web/viewer.html?file=../25.pdf#page=177}{sites 176--177} and \href{https://grothendieck.umontpellier.fr/web/viewer.html?file=../25.pdf#page=179}{site 178}.

The proof of 16.9 constructs the isomorphism functor between the standard pair \((E_o,Q_o)\) and \((E,Q)\), represented in the text by \(P\). Site 178 places it in the ambient space associated with \(\underline{\mathrm{Hom}}(E_o,E)\) and an open isomorphism scheme; the latter display actually uses \(E_1\), which remains in the transcription. Site 179 introduces \(\underline Q(E)=W(\underline{\mathrm{Sym}}^2(E))\), the section associated with a quadratic form, and the open part \(\underline Q(E)^*\) defined by the corrected discriminant. Transport of structure yields the displayed orbit morphism from \(\underline{\mathrm{Isom}}(E_o,E)\) to this open part. Proposition 16.11 states faithful flatness of that morphism, and \(P\) is described as its inverse image along the section corresponding to \(Q\). The proof on site 180 reduces to an algebraically closed field, invokes generic flatness, and uses the action of \(Gl(E)\) to translate a suitable open subset. The marginal instruction asks for a lemma about morphisms of homogeneous spaces; the compressed word and the edge-clipped beginning of the last word remain qualified. Lemma 16.11.1 supplies the standard-form step. The proof cites Bourbaki for one part and treats the remaining part by a direct-sum decomposition and Corollary 16.8. See \href{https://grothendieck.umontpellier.fr/web/viewer.html?file=../25.pdf#page=179}{site 178}, \href{https://grothendieck.umontpellier.fr/web/viewer.html?file=../25.pdf#page=180}{site 179}, and \href{https://grothendieck.umontpellier.fr/web/viewer.html?file=../25.pdf#page=181}{sites 180--181}.

Site 181 completes the proof using Lemma 16.11.2, whose inserted title and smoothness/standard-form comparison are enclosed in a handwritten oval. The oval's upper boundary crosses the preceding typed line; that line remains part of the text. The argument explicitly refers back to 16.2 and 16.3. A following note proposes moving 16.11.2 to the beginning and interchanging the earlier two items, so that the discussion independent of the corrected discriminant would be grouped together. The same note asks the recipient about prior knowledge, attribution, and established terminology. Its continuation on site 182 says that Bourbaki is not at hand for checking and proposes a terminological revision involving ``déterminant'' and ``discriminant''. The guide reports that proposal as part of the correspondence embedded in the manuscript; the author's uncertainty about attribution and terminology remains visible. See \href{https://grothendieck.umontpellier.fr/web/viewer.html?file=../25.pdf#page=182}{site 181} and \href{https://grothendieck.umontpellier.fr/web/viewer.html?file=../25.pdf#page=183}{site 182}.

Remark 16.12 on site 182 introduces \(\underline O(n)\) as the stabilizer of the standard form and \(\underline O(Q)\) as the relative stabilizer. It returns to the morphism of 16.11, the right action of the orthogonal group, and the principal homogeneous bundle associated with a form. The inserted pair of point labels has overlapping \(Q/u\) readings, while the following equation uses \(u'=uv\); those letters have been preserved with their respective uncertainties. The text refers to a planned fuller treatment in EGA V or VI and adds ``voir en attendant SGAD IV'' by hand. On site 183 the construction \((E,Q)\rightsquigarrow P\) is connected with an equivalence of the stated fibered categories through formal arguments and flat descent; the chapter-reference cell after V is overworked and remains read as \(1/I\). Paragraph 16.13 then develops the étale-trivialization discussion under its rank and characteristic restriction. It invokes the lemma of Hensel for the added condition (iii ter) and cites SGAD XXIV for the finite étale refinement. These are explicit dependencies named by the manuscript. See \href{https://grothendieck.umontpellier.fr/web/viewer.html?file=../25.pdf#page=183}{site 182} and \href{https://grothendieck.umontpellier.fr/web/viewer.html?file=../25.pdf#page=184}{site 183}.

The handwritten addition 16.14 on site 183 supplies the rank-one example \(at^2\) over an imperfect field of characteristic two, with the displayed \(a\in k-k^2\). Site 184 continues it by referring to \(k(a^{1/2})\), then returns to a finite, locally free, surjective morphism \(S'\longrightarrow S\). Its added group expression writes a quotient of \(\underline O(n)_S\) by \(\underline{SO}(n)_S\) beside the literal \(\mu_2\,S\); the transcription keeps the separate baseline S and the mixed typed/handwritten layers. The text describes the reduction of the structural group to \(\underline{SO}(n)\), returns to the earlier local-isotriviality claim, and ends with the local-base formulation using (iii bis). The lower half of the final supplied page has no further text apart from the archive folio. See \href{https://grothendieck.umontpellier.fr/web/viewer.html?file=../25.pdf#page=184}{site 183} and \href{https://grothendieck.umontpellier.fr/web/viewer.html?file=../25.pdf#page=185}{site 184}.
\par\typeout{EGA-SOURCE-END guide-iv17-16 PDF-PAGE \thepage}

\par\typeout{EGA-SOURCE-BEGIN guide-iv20-opening PDF-PAGE \thepage}
\pdfbookmark[1]{IV 20: opening prospectus}{source:guide-iv20-opening}
\subsection*{IV 20: the opening contents and Commentaires (cote 25, sites 2--3)}
The opening contents list sixteen proposed sections, from \emph{Notations} to \emph{Traductions en langage de systèmes linéaires}. In the \emph{Commentaires}, Grothendieck describes the present \emph{Etat 0} as an assembly of detailed drafts for a later redaction, and expressly calls for a new plan. He names the current nos.\ 11, 12, 14 and 15 as material to be moved earlier. The continuation on site 3 places proposition 10.3 among results needing an earlier home and mentions further material on \emph{applications birationelles}. He proposes splitting the current paragraph 20, with the developments on grassmannians making one of the resulting parts. The projected no.\ 16 remains unwritten at the moment of these comments. These are changes of organization stated on the pages themselves; the edition preserves the successive headings and references. The doubt retained in the upper-right pencil mark of site 3 concerns a folio annotation; it does not supply a new section number. Sources: \href{https://grothendieck.umontpellier.fr/web/viewer.html?file=../25.pdf\#page=3}{cote 25, site 2} and \href{https://grothendieck.umontpellier.fr/web/viewer.html?file=../25.pdf\#page=4}{site 3}.
\par\typeout{EGA-SOURCE-END guide-iv20-opening PDF-PAGE \thepage}

\par\typeout{EGA-SOURCE-BEGIN guide-iv20-1 PDF-PAGE \thepage}
\pdfbookmark[1]{IV 20.1}{source:guide-iv20-1}
\subsection*{IV 20.1: Préliminaires et notations (cote 25, sites 4--8)}
The section introduces the projective bundle $P$ attached to a locally free Module of finite type, its dual projective bundle, the \emph{schéma des hyperplans}, and the hyperplane attached to an invertible quotient of the dual Module. The displayed pairing and its transpose lead to the defining ideal; a second description uses a local section, and a third identifies the hyperplane with a projective bundle. Grothendieck expressly prefers beginning with this third description (site 5). His references to paragraph 17 concern smoothness, and those to paragraph 19 concern \emph{diviseurs relatifs} and \emph{intersection complète transversale}; Chapter III is invoked for a proposed determination of the relative-divisor functor. The \emph{schéma d'incidence} and \emph{hyperplan universel} organize the inverse images $Y_{\xi}$ and the corresponding Modules in a diagram on site 7. The final paragraph on site 8 extends the term \emph{section hyperplane} from an immersion to the stated general morphism. Several visible slips or unsettled signs merit attention: site 4, formula (xx), prints $\underline{D}_P(1)$, followed by $\underline{L}_P(1)$ in the next sentence; the letter beneath the inverse-image discussion on site 7 has overprinted M/H traces, retained as a doubt. These forms have not been repaired from the mathematics. The handwriting also replaces \emph{linéaires} by \emph{coniques} in the opening title, adds \emph{fermé}, and revises the wording of the universal-hyperplane paragraph. Sources: \href{https://grothendieck.umontpellier.fr/web/viewer.html?file=../25.pdf\#page=5}{sites 4--8}, especially sites 5--7 for the reorganization and universal construction.
\par\typeout{EGA-SOURCE-END guide-iv20-1 PDF-PAGE \thepage}

\par\typeout{EGA-SOURCE-BEGIN guide-iv20-2 PDF-PAGE \thepage}
\pdfbookmark[1]{IV 20.2}{source:guide-iv20-2}
\subsection*{IV 20.2: Etude de la section hyperplane générique---propriétés locales (cote 25, sites 8--15)}
The section studies the generic hyperplane through the incidence construction of no.\ 1. Propositions 2.1--2.6 concern irreducibility, empty inverse images, the induced correspondence of irreducible components, and dimension; proposition 2.7 treats \emph{cycles premiers associés}. The proof refers to paragraph 3 for associated cycles under the stated flat morphism, then explicitly shortens its own route through the universal incidence family (site 10). Proposition 2.8 and its corollaries use paragraph 6 for the local properties \emph{coprof}, $(S_k)$, Cohen--Macaulay, reducedness and integrality. Corollary 2.11 adds regularity, $(R_k)$ and normality. The discussion then turns to smoothness and the corresponding geometric properties in 2.12--2.15, using the \emph{critère jacobien de lissité}, followed by \emph{multiplicité radicielle} and separability in 2.16--2.17. The handwriting changes hypotheses and the organization: a canceled transition on site 11 reappears on site 12; the inserted hypothesis preceding 2.12 and the margin beside 2.14 contain a ramifié qualification whose initial letters remain doubtful; site 14 proposes presenting the practical corollary before its supporting result; Remarques 2.18 revisit the immersion assumption and characteristic zero. The checked pages retain several mathematical slips or gaps: the index is absent from $k(\ )$ on site 13, the argument there prints $p\leq n$ before its separate $p=n$ discussion, and both occurrences in the non-embeddedness assertion of 2.16 on site 14 read $Z_i$ without an added generic-point index. The guide records these forms without repairing their hypotheses or formulas. Sources: \href{https://grothendieck.umontpellier.fr/web/viewer.html?file=../25.pdf\#page=9}{cote 25, sites 8--15}; internal references to published EGA IV, paragraphs 3 and 6, are those explicitly written in the prenotes.
\par\typeout{EGA-SOURCE-END guide-iv20-2 PDF-PAGE \thepage}

\par\typeout{EGA-SOURCE-BEGIN guide-iv20-3 PDF-PAGE \thepage}
\pdfbookmark[1]{IV 20.3}{source:guide-iv20-3}
\subsection*{IV 20.3: Section hyperplane générique---irréductibilité et connexité géomét. (cote 25, sites 15--18)}
The Bertini--Zariski discussion turns from the local properties of no.\ 2 to geometric irreducibility and connectedness of the generic section. It writes the function field of the incidence section in affine coordinates, reducing the argument to the lemma attributed to Zariski on site 17; the explicit birational description is itself proposed as a corollary to 2.1. The handwritten revision of that lemma replaces the original algebraicity qualification by a transcendence-degree qualification and deletes the claim equating radicial multiplicities, while retaining the assertion about primary extensions. The surrounding argument uses the generic incidence construction of no.\ 1 and the separability discussion of no.\ 2. The connectedness corollary invokes paragraph 6 and the theorem named ``Hartshorme'' in the typescript, joining irreducible components through their intersections. Its script family of closed subsets remains typographically doubtful in the checked text. Several changes matter to the reading: the theorem on site 15 loses its geometrically integral alternative; a long canceled proof continues through the first part of site 16; the new proof begins at ``Soit K/k'' there; site 17 adds the attribution ``(Zariski)'', modifies the immersion hypothesis of corollary 3.2, and inserts algebraic closedness in corollary 3.3. The author leaves Dieudonné the choice of a proof to copy from the literature. The qualification above the canceled immersion wording still has an uncertain initial group. The guide preserves the written argument and its omissions: the canceled formula on site 16 has a blank in $k(\ )$, and site 18 prints $Y_\eta=X_\eta$ while completing the connectedness proof. Sources: \href{https://grothendieck.umontpellier.fr/web/viewer.html?file=../25.pdf\#page=16}{site 15}, \href{https://grothendieck.umontpellier.fr/web/viewer.html?file=../25.pdf\#page=17}{site 16}, \href{https://grothendieck.umontpellier.fr/web/viewer.html?file=../25.pdf\#page=18}{site 17}, and \href{https://grothendieck.umontpellier.fr/web/viewer.html?file=../25.pdf\#page=19}{site 18}; the paragraph-6 reference is the one given in the prenotes.
\par\typeout{EGA-SOURCE-END guide-iv20-3 PDF-PAGE \thepage}

\par\typeout{EGA-SOURCE-BEGIN guide-iv20-4 PDF-PAGE \thepage}
\pdfbookmark[1]{IV 20.4}{source:guide-iv20-4}
\subsection*{IV 20.4: Etude d'une section hyperplane variable---sections ``assez générales'' (cote 25, sites 18--21)}
This section passes from the generic section to constructible loci in the dual projective space, using paragraph 9 and the finite-presentation incidence morphism. The local properties collected in proposition 4.2 include the printed ``coppof'', $(S_k)$, Cohen--Macaulay modules, associated prime cycles without embedded components, and separability; geometric integrality is crossed out in that list. Proposition 4.3 makes the corresponding list for schemes, including smoothness, geometric $(R_k)$, normality, irreducibility and integrality, with the last two labels visibly revised and linked by a pen arrow. The proof points back to nos.\ 2--3, while a ``Remords'' on site 19 broadens the proposed exceptional set. Corollary 4.4 introduces the connectedness property $(C_m)$ through deletion of closed subsets of bounded dimension. The next page explicitly identifies a lapse in its statement and directs the reader to measure the dimensions of the images in $P$ or $H_\xi$ in the general discussion; the corrected readers retain the author's stated lapse and his accompanying repair instruction rather than silently rewriting the corollary. The same page invokes AQT for constructibility, calls for an addition to paragraph 9, and asks for a corresponding generalization of 3.3. The editorial sequence is unusually visible: site 18 cancels a promise to treat openness by paragraph 12; site 21 cancels the proposed paragraph 4.5 and proposition 4.6, while its margin sends 4.5 to no.\ 5.1; site 22 then calls that paragraph's cancellation an error. The draft proposition 4.6 ends incompletely on site 21. Other printed forms retained for the reader include the inequality $\operatorname{coppof}(M)\leq n$ on site 19 and the indexed $E_\xi$ at the end of site 20. Sources: \href{https://grothendieck.umontpellier.fr/web/viewer.html?file=../25.pdf\#page=19}{sites 18--21}, with the restoration instruction on \href{https://grothendieck.umontpellier.fr/web/viewer.html?file=../25.pdf\#page=23}{site 22}; the explicit references to published EGA IV are paragraphs 9 and 12.
\par\typeout{EGA-SOURCE-END guide-iv20-4 PDF-PAGE \thepage}

\par\typeout{EGA-SOURCE-BEGIN guide-iv20-5 PDF-PAGE \thepage}
\pdfbookmark[1]{IV 20.5}{source:guide-iv20-5}
\subsection*{IV 20.5: Théorèmes du type Seidenberg (cote 25, sites 22--25, with the restored paragraph from site 21)}
The opening instruction restores the preceding paragraph about openness, whose margin directs it to no.\ 5.1. The section then uses the incidence divisor and its defining section $\phi$, expresses the inverse-image module as a cokernel, and defines $F$-regularity. Proposition 5.3 treats the regularity locus $U$ and relative flatness; corollary 5.4 gives the corresponding locus $V$ in the dual projective space. The pen loop on site 23 moves the finite-presentation and flatness hypothesis into part a), and the handwritten additions to 5.4 distinguish its constructibility and openness assertions. Remark 5.5 discusses associated cycles and supports; theorem 5.6 and corollary 5.7 invoke the lists in 12.1.1 and 12.1.4. The text expressly uses 11.3 for the flatness argument, mentions EGA III without a completed number in the locally noetherian case, and refers to paragraph 3 and nos.\ 2.3 and 2.6 for the associated-cycle argument. Site 24 asks that the final exposition replace references to lists in another fascicule by an explicit list and identify Seidenberg's geometrically normal case. The crossed commentary on site 25 explains the proposed division of the preceding material and compares the route through 4.6 with the ``délicat 15.2.2''. It also states that the author has not read Seidenberg's paper. The checked text retains mathematical slips visible in the formulation: site 22 speaks of an isomorphism of $F$ with $\underline{O}_X$ after discussing an invertible $L$ and later writes $\mathrm{id}_F\otimes\phi$ before the displayed source $F\otimes L^{-1}$; site 23 prints $s\in\check P$ in 5.3 b); site 24 concludes with flatness relative to $S$ in 5.6. The canceled handwritten paragraph on site 23 reads ``$k$ un corps'', correcting the draft transcription's ``$A$ complet''; its comparison sign remains uncertain. The canceled Remords paragraph on site 27 returns to the same proposed reorganization: it separates openness proved through a valuative criterion from the additional flatness conclusion using 15.2.2, and proposes a new 4.8. Sources: \href{https://grothendieck.umontpellier.fr/web/viewer.html?file=../25.pdf\#page=22}{site 21}, \href{https://grothendieck.umontpellier.fr/web/viewer.html?file=../25.pdf\#page=23}{sites 22--25}; the manuscript explicitly cites 11.3, 12.1.1, 12.1.4 and 15.2.2. The published numbering needs care: EGA IV 12.1.4 is a remark about fiber dimension, whereas the published list of geometric fiber properties occurs in 12.1.6. The manuscript number is preserved. See \href{https://www.numdam.org/articles/10.1007/BF02684343/}{EGA IV, third part, 12.1.4 and 12.1.6, pp.\ 177--178}.
\par\typeout{EGA-SOURCE-END guide-iv20-5 PDF-PAGE \thepage}

\par\typeout{EGA-SOURCE-BEGIN guide-iv20-6 PDF-PAGE \thepage}
\pdfbookmark[1]{IV 20.6}{source:guide-iv20-6}
\subsection*{IV 20.6: Connexion d'une section hyperplane quelconque (cote 25, sites 25--26)}
The short connectedness section combines the geometric connectedness of the generic hyperplane section from 3.3 with the Zariski connectedness theorem in the form cited as 15.6.3. Its argument uses the proper incidence family, dominance of the dual projective space by its irreducible components, and the geometric unibranchness of the regular parameter space, then applies the cited result to all fibers. The author adds properness of $X$ in the margin of site 25 and strengthens the dimension comparison by a handwritten lower bar. The final paragraph on site 26 asks for a generalization in the sense of 4.4, expressly includes sections of excess dimension, and proposes collecting all the connectedness material, including 3.3 and 4.4, in one number. The next proposed application heading is canceled on site 26 and reappears as 7 on site 27. Numbering is part of the unfinished revision: the handwritten proposition and the two references on 26 read 6.1, while the overwritten section numeral on 25 remains uncertain between 6 and 5 and its initial capital has overlapping C/E traces. These are retained rather than reconciled by renumbering. The first connecting pen stroke in the properness paragraph has no visible arrowhead, while the two later occurrences do; the checked source preserves that graphical evidence. Sources: \href{https://grothendieck.umontpellier.fr/web/viewer.html?file=../25.pdf\#page=26}{site 25} and \href{https://grothendieck.umontpellier.fr/web/viewer.html?file=../25.pdf\#page=27}{site 26}; the manuscript cites EGA IV 15.6.3. The precise concordance of that citation with the published numbering remains unestablished here.
\par\typeout{EGA-SOURCE-END guide-iv20-6 PDF-PAGE \thepage}

\par\typeout{EGA-SOURCE-BEGIN guide-iv20-7 PDF-PAGE \thepage}
\pdfbookmark[1]{IV 20.7}{source:guide-iv20-7}
\subsection*{IV 20.7: Application to hyperplane sections and multisections of specified type (cote 25, sites 27--30)}
This section turns the preceding open subsets of the dual projective space into families of hyperplane sections and then into multisections. The discussion moves from rational points over an infinite field to sections near a chosen base point, and then to finite separable residue-field extensions and an étale change of base. The author redirects the flatness argument from the old references 4.7--4.8 to the revised number 5 on site 28. A canceled preliminary passage on that page gives way to the remords on site 29: the singular locus for a property P is allowed to have a specified dimension, and the resulting formulations 7.2--7.3 track its dimension after taking a hyperplane section. The margin adds another case to 4.3 through published EGA IV 12.2.1, with the final parenthesized reference still doubtful. The two expanded subscripts on site 29 visibly read $\xi,\xi(s)$, while the intervening G has only $\xi(s)$; the neighborhood is printed as a neighborhood of s ``dans E''. Those forms remain in the text. Proposition 7.4 and remarks 7.5 develop the construction of étale multisections, including the successive reduction of fiber dimension; the final marginal note proposes avoiding an ensemble whose faint adjective remains doubtful. The canceled discussion across sites 29--30 invokes EGA III 7 and Stein factorization for geometric connected components. Its replacement margin cites EGA IV 12.2.1, with the first Roman numeral unresolved and the second legibly (xi). The earlier canceled Remords on site 27 separates openness by the valuative criterion from the flatness conclusion using EGA IV 15.2.2; it records a proposed reorganization of 4.6--4.8, rather than a completed renumbering. Sources: \href{https://grothendieck.umontpellier.fr/web/viewer.html?file=../25.pdf\#page=28}{site 27}, \href{https://grothendieck.umontpellier.fr/web/viewer.html?file=../25.pdf\#page=29}{site 28}, \href{https://grothendieck.umontpellier.fr/web/viewer.html?file=../25.pdf\#page=30}{site 29}, and \href{https://grothendieck.umontpellier.fr/web/viewer.html?file=../25.pdf\#page=31}{site 30}.
\par\typeout{EGA-SOURCE-END guide-iv20-7 PDF-PAGE \thepage}

\par\typeout{EGA-SOURCE-BEGIN guide-iv20-8 PDF-PAGE \thepage}
\pdfbookmark[1]{IV 20.8}{source:guide-iv20-8}
\subsection*{IV 20.8: Dimension de l'ensemble des hyperplans exceptionnels (cote 25, sites 31--54)}
This long, repeatedly revised section studies the codimension of the exceptional hyperplanes for the properties developed in sections 2--7, with linear pencils as its stated motivation. The opening distinguishes the constructible exceptional set from its closure and asks whether these refinements belong here or beside the earlier results; the first version of 8.2 is canceled on site 31. Its replacement on 32 treats excess dimension and then the regularity of the section $\phi_\xi$ on a coherent Module $F$, using associated prime cycles and the incidence family. Here the first assertion literally says codimension 2, while the later Module formulation has an explicit comparison sign before 2; the first statement is not supplied with one. The same discussion changes among $Z$, $T$, and their irreducible components. Site 33 passes through coprofondeur, the Cohen--Macaulay property, and $(S_k)$: the added sentence near the top has $(S_{k-1})$, whereas the final handwritten continuation has $(S_k)$. The canceled intermediate argument and its replacement must be read with that distinction of written indices in view. Proposition 8.4 and its dimensional lemma, introduced as a remords to paragraph 5, prepare a more precise description of the exceptional hyperplanes and then the criterion for $(S_k)$ in 8.5--8.6. The author cites the dimension theory of $0_{\mathrm{IV}}$ and published EGA IV 5.7.4. Several literal forms require attention: site 34 says $T_i$ is a component ``de X'', gives the lemma the inserted hypothesis $\operatorname{dodim}(Y,X)\geq k$, and cancels $\dim f(X_j)$ while leaving the following equals-zero expression; its last line has $f(T)$ without an index. The corollary on 35 is replaced rather than merely augmented: the new introductory sentence cites 8.4, 5.7.4, and the end of 8.2, and the partly clipped handwritten hypothesis remains uncertain. Site 36 retains competing readings in $n+k+1$, $f(T_j)$, and the $F/G$ expression; its necessity argument literally calls $\phi_\xi$ $X_{k(\xi)}$-regular amid formulas about $F$ and $G_\xi$. The next part, beginning with the revised number 8.6 on 37, gives the differential description behind the exceptional locus. It uses the exact sequences on projective space, the conormal Module, and the singular and supersingular zero constructions of paragraph 16. The projective conormal description and the diagrams on 39--42 connect ordinary singular points, ramification, tangent hyperplanes, and the stated orthogonality. The author's canceled preliminary argument on 39 gives $M$ rank $n$, writes $\underline{O}_X^{\,m}$, and uses $n-m$ for the fiber dimension; those three expressions are retained together. Other written notation is likewise left visible: the $r-1$ and $d-1$ occurrences on 38, the change from $\psi$ on 39 to $\mu$ on 40, the term $\Omega_{X/Y}$ in 40's diagram, and the $X/Y$ subscript and exponent $-2$ in 42's lower-left object. The hypotheses and numbering also evolve: 8.7 changes the point label from $x$ to $y$, 8.8 receives the immersion hypothesis and the reminder that $X$ remains smooth, and the margin of 43 restricts the role of immersion to the equivalence with clause (i ter). The dimension sentence on 43 has the handwritten alternative $\underline{Y}^{\mathrm{sup\ sing}}\neq\emptyset$ followed by the revised equality $r-2$; the later inserted alternative on 45 has an equals-empty-set sign and belongs immediately before the dimension statement there. The edition preserves both signs. The proof on 44--45 works with the tangent map and its image and with a determinant section; the first handwritten morphism label on 44 remains uncertain, and 45 retains both the earlier $r$ and the later $r-1$ dimension expressions. Sites 46--49 return to regularity, geometric normality, irreducibility, and the extent of the required nonramification hypotheses. The term ``réglé'' and successive versions of the exceptional-set criterion are central here. Long canceled formulations on 47--48 remain in the diplomatic text and apparatus; the repeated call marks on 47 join the top addition to the sentence about what $T$ contains, while those on 48 place the addition beginning ``grâce au fait'' after the codimension bound. That latter addition visibly uses $\underline{X}'^{\mathrm{sing}}$, and the text retains it. The correction note on 49 keeps its two primed $X'_i$ expressions and a doubtful $X/Y$ in the final image formula. The remords on 50 weakens the proposed nonramification assumption; 8.16 and 8.17 then combine geometric normality, depth, connectedness, and irreducibility. The codimension character in the remords remains an unresolved overprint. The quadric example contains the handwritten ``en effet tous'', giving the literal wording about tangent hyperplanes being reducible; on 51 the author immediately records his alarm that the quadric does not qualify under his own use of ``réglé'' and asks for another word. This reservation is part of the argument's unfinished state. Proposition 8.18 on 51 has an inserted dimension $d\geq2$, still says dimension 1 in its main assertion, and changes the dimensions in clauses (a) and (b) to $d-1$; clause (a) retains an unchecked $P(k)$. The top handwritten addition, placed before ``Nous allons'' by matching call marks, uses an underlined $Y^{\mathrm{sing}}$ without a $\xi$ index. Its final short mark remains doubtful. The canceled proof below it gives way to 52's proof through the normal case, the open subset $U$, and normalization; the handwritten note compares 8.18 with 8.16 and explains the interest in allowing $d\geq2$. The closing discussion on 53--54 studies hyperplanes containing a prescribed linear subspace through the projection $P-L^o\to Q$ and invokes the earlier Bertini--Zariski argument of 3.1. Site 53 changes from $Y_\xi$ to $X'_\xi$ and $X_\xi$, then literally ends the equivalence with unindexed $X$; its proposed converse also retains the semicolon and doubtful numeral after $\dim f(X)$. The generic-fiber condition is added below that page and developed in the transcendence-degree discussion on 54. The final dimension estimates, the condition at the finitely many points of $T$, and the suggestion to move these ``Bertiniesques'' developments to the following section close the passage. No index, inequality, hypothesis, or diagram has been repaired from this commentary. Sources: \href{https://grothendieck.umontpellier.fr/web/viewer.html?file=../25.pdf\#page=32}{sites 31--36}, \href{https://grothendieck.umontpellier.fr/web/viewer.html?file=../25.pdf\#page=38}{sites 37--42}, \href{https://grothendieck.umontpellier.fr/web/viewer.html?file=../25.pdf\#page=44}{sites 43--49}, \href{https://grothendieck.umontpellier.fr/web/viewer.html?file=../25.pdf\#page=51}{sites 50--52}, and \href{https://grothendieck.umontpellier.fr/web/viewer.html?file=../25.pdf\#page=54}{sites 53--54}.
\par\typeout{EGA-SOURCE-END guide-iv20-8 PDF-PAGE \thepage}

\par\typeout{EGA-SOURCE-BEGIN guide-iv20-9 PDF-PAGE \thepage}
\pdfbookmark[1]{IV 20.9}{source:guide-iv20-9}
\paragraph{IV 20.9: changing the projective immersion (cote 25, sites 55--57).}
Section 9 occupies sites 55--57. Paragraph 9.1 introduces the immersion $u_n:P\to P(n)$ associated with the degree-$n$ symmetric power and describes replacing $f$ by $u_nf$. Its examples refer back to 8.12.b and to singular hyperplane sections of a quadric; the latter passage has a handwritten replacement by ``dans $\mathbf P^3$''. The displayed direct-image expression has no visible equality before the symmetric power, and the edition preserves that gap. Proposition 9.2 is stated over $S=\operatorname{Spec}(k)$ for smooth $X$ and an unramified $f:X\to P$, with $n\geq2$. It records the conditions cited as equivalent in 8.8 and describes the singular points of the corresponding sections outside an exceptional locus of codimension at least two. The surrounding revision replaces the original single-point wording by a finite-number formulation and retains a separate assertion for an immersion. The note below it explicitly calls for a formulation of 8.8 that accommodates the unramified case, referring to condition (ii bis).

Proposition 9.3 and its proof run from site 56 into site 57. The opening has a typed insertion imposing geometric irreducibility of $X$ and a handwritten $n\geq2$; the dimension assumption concerns $f(X)$. Its three clauses concern sections through a given point, sections tangent there to a specified linear subvariety, and sections tangent to $X$. The geometric-irreducibility conclusions are proved in the draft by choosing polynomials $P,Q,R$ and examining $Q/P$ and $R/P$ in the function field. The proof writes out the degree-two choices and the exceptional planar discussion, followed by a marginal set of squared-coordinate choices for $\dim L\geq3$. That marginal note is linked to the body by matching oblique marks. Several overprinted coordinates and the handwritten point symbol remain doubtful in the transcription. Site 57 then proposes combining the argument with a reference beginning 8.18, whose final overprinted character is unresolved, under the unfinished heading Corollaire 9.4. A further appeal to 9.2 introduces the proposed ``cas excellent'' statement numbered 9.5, with a handwritten Th\'eor\`eme over a partly obscured typed title and a final note about a fundamental-group-like symbol and singular points. The edition preserves these drafting instructions and their uncertainty; it supplies no completion to either unfinished statement.
\par\typeout{EGA-SOURCE-END guide-iv20-9 PDF-PAGE \thepage}

\par\typeout{EGA-SOURCE-BEGIN guide-iv20-10 PDF-PAGE \thepage}
\pdfbookmark[1]{IV 20.10}{source:guide-iv20-10}
\paragraph{IV 20.10: pencils of hyperplane sections and blown-up varieties (cote 25, sites 58--62, opening).}
Paragraph 10.1 relates the constructible exceptional set for a property of hyperplane sections to lines in the dual projective space, and defines the pencil associated with such a line. The discussion uses the results of sections 8 and 9, cites 8.2 for the additional regularity condition, and refers to 7.1 and paragraph 5 in its passage from a field to a neighborhood on a general base. The handwritten reference at site 58 reads 12 over the canceled word ``suivant''. Paragraph 10.2 begins at site 59 with the polar linear subvariety $L^o$ of codimension two and the fiber product $T=X\mathbin{\mathrm x}_P L^o$. It describes $T$ through the homomorphism from the rank-two submodule and through its ideal of zeros. The regularity hypothesis is expressed using a totally ordered local basis, then an isomorphism with $J/J^2$ together with a regular immersion. Th\'eor\`eme 10.2 states the unique $X$-isomorphism between the pencil and the blow-up of $X$ along $T$. The handwritten uniqueness phrase and the tensor insertion in the preceding condition are retained with their actual attachment points.

Site 60 discusses how to organize that theorem and the proposed statement 9.5, asks for the regularity requirement in the recapitulative statements, and describes the intended fibrations over $\mathbf P^1$. Its corrected paragraph reference is again 12; the attached note has an unresolved first character in S.3./5.3. The words ``sont lisses sont au plus'' are printed in that order and remain in the edition. Site 61 gives the more general Proposition 10.3, formulated for a quasi-coherent Module, its homomorphism to $\underline O_X$, and the induced construction over $P(\underline G)$. The following note calls its first assertion nearly immediate and leaves the second as an exercise for Dieudonn\'e. Two large crossing diagonals traverse this proposition and note, while the preceding paragraph suggests moving the material to the discussion of regular immersions. Their editorial function remains unresolved, so the entire block is retained with the graphic marks recorded. Four number groups on this page clearly show handwritten 10 over typed 8; the ambiguous connector in its displayed sequence is preserved without completing it. The final paragraph combines 10.2 with 9.2 in dimension at most one and leads to Corollaire 10.4 for a smooth geometrically connected curve in projective space over an infinite field. The corollary continues at the opening of site 62: the degree is revised by hand to $d$, and the fiber description contains $d-2$ copies of the indicated spectrum together with $\operatorname{Spec}k'[t]/(t^2)$. Its following explanation uses the term ``quadratique'' for the ramification. The old canceled fiber description and its uncertain coefficient remain in the apparatus.
\par\typeout{EGA-SOURCE-END guide-iv20-10 PDF-PAGE \thepage}

\par\typeout{EGA-SOURCE-BEGIN guide-iv20-11 PDF-PAGE \thepage}
\pdfbookmark[1]{IV 20.11}{source:guide-iv20-11}
\paragraph{IV 20.11: Grassmannians and linear subvarieties (cote 25, sites 62--65).}
Section 11 begins partway down site 62 under ``Grassmaniennes'', with the old following heading ``Ratrapages'' canceled. The opening explains the need for linear subvarieties beyond the dimensions previously used and for a review of the earlier sections in this enlarged setting. It defines the functor of locally free rank-$n$ quotients of $\underline E$ and constructs its morphism to the rank-one Grassmannian, identified in the display with the projective bundle of an exterior power. The proof refers explicitly to the S\'eminaire Cartan and continues on site 63 with a closed-subscheme realization and the stated separation, quasi-compactness, projectivity and finite-presentation properties. The source's printed ``immersion ferm\'ees'', its internally spaced words, and the unresolved base letter at the start of the definition are retained. For a locally free Module, site 63 gives the affine charts indexed by choices of $n$ basis vectors, with relative dimension $n(r-n)$, and describes their corresponding subfunctors through a chosen exact sequence. Its displayed vector-bundle expression replaces a canceled tensor expression by a Hom expression. The consequences written there include local lifting of a rational point to a section and the existence of rational points in nonempty opens over an infinite field.

Sites 64--65 interpret the Grassmannian in terms of linear subvarieties $P(\underline G)$ of $P(\underline E)$, with a handwritten qualification that the general terminology imposes no fixed rank. The intrinsic characterization by smoothness, relative dimension and projective degree one is explicitly deferred to a later chapter. The polarity discussion writes a canonical comparison of the rank-$n$ and rank-$(r-n)$ Grassmannians and interprets the corresponding dual linear subvariety as the hyperplanes containing the original one. The visible unaccented $E$ in ``soit E son dual'' and the neighboring handwritten V-shaped dual accents remain as written. Site 65 records the containment criterion through the comparison of two locally free quotients, with the directions of the displayed morphisms preserved. The author then calls for a minimum collection of sorites whose full list will depend on the completed preceding material. A canceled section-12 heading is followed by the handwritten addition introducing the quotient functor with no prescribed rank and representing it by the coproduct over $n\geq0$. Its small index is partly uncertain. The concluding note expressly allows the rank, and hence relative dimension, to vary over a disconnected base. Site 66 begins the separate section on generalizing earlier results to linear sections, under the revised title ``Généralisation de résultats antérieurs aux sections linéaires''.
\par\typeout{EGA-SOURCE-END guide-iv20-11 PDF-PAGE \thepage}

\par\typeout{EGA-SOURCE-BEGIN guide-iv20-11-notes PDF-PAGE \thepage}
\pdfbookmark[1]{IV 20.11: final-page observations}{source:guide-iv20-11-notes}
% The supplementary observations formerly printed here are consolidated into guide_parts/20_11.tex.
% The complete prior text and clause mapping are retained in records/guide_copyedit/ and records/GUIDE_IV20_11_CLAUSE_MAP.tsv.
\par\typeout{EGA-SOURCE-END guide-iv20-11-notes PDF-PAGE \thepage}

\par\typeout{EGA-SOURCE-BEGIN guide-iv20-12 PDF-PAGE \thepage}
\pdfbookmark[1]{IV 20.12}{source:guide-iv20-12}
\paragraph{IV 20.12: Généralisation de résultats antérieurs aux sections linéaires (cote 25, sites 66--74).}
Grothendieck extends the notation and the working constructions of the preceding hyperplane sections to linear sections of codimension $m$. Site 66 distinguishes the indices for dimension and codimension, introduces the Grassmannian parameter space and its incidence section, and relates the two transposed homomorphisms defining that section. At site 67 he explicitly criticizes the choice of a new letter $Y$ for the incidence construction and proposes $\underline{X}^{(m)}$; his displayed comparison of iterated exponents is accompanied by the words ``à peu de choses près''. The identification with a Grassmannian over $X$ leads into a detailed list of changes to the statements of sections 2--10 (sites 68--70), with several instructions addressed directly to Dieudonné. The added condition on dimensions, the repeated use of $F$-regularity, and the appeal to paragraph 19 belong to that list. A canceled initial objection at site 71 gives way to an extended discussion of the scope and editorial cost of the proposed generalization. The singular and supersingular loci $V^1$, $V''$, and $V^2$ are then related to the complements to paragraph 16 and to the filtration of a Grassmannian by dimensions of intersection with a fixed linear subspace (sites 71--73). His qualifications ``sauf erreur'' and ``sauf erreur de calcul'' accompany this discussion and its dimension expression; they remain visible in the text. Site 74 turns to the projective family $Q(C)$ of linear spaces through a fixed $C$, recalls 10.2 with a retouched following reference, and adds the previously omitted nonramification and Veronese-composition requirements to the proposed statement. The final handwritten continuation retains $Z$ in its last ``assez général'' phrase. Other notation requiring attention includes the unaccented $E$ followed by ``soit F'' in the incidence construction at site 66, the canceled material after the Grassmannian subscript $m$ on that page, and the baseline typed $v$ after $G_T$. These are readings of the scans, and are left in place; the apparatus records the written changes without supplying a mathematical revision.
\par\typeout{EGA-SOURCE-END guide-iv20-12 PDF-PAGE \thepage}

\par\typeout{EGA-SOURCE-BEGIN guide-iv20-13 PDF-PAGE \thepage}
\pdfbookmark[1]{IV 20.13}{source:guide-iv20-13}
\paragraph{IV 20.13: Morphismes élémentaires et théorème de M. Artin (cote 25, sites 75--83).}
Definition 13.1 introduces elementary morphisms by removing an étale closed subscheme from a smooth projective family of geometrically connected curves, and calls composites of elementary morphisms poly-elementary. Theorem 13.2, with Artin's name and the hypotheses ``géom. irréd.'' and ``parfait infini'' added above the line at site 75, gives poly-elementary open neighborhoods of a smooth point. Its induction begins with the normal projective model of a function field and explicitly refers to Chapter II, paragraph 7; the text also explains its use of the composite of an empty family. The reduction to a normal projective scheme, the singular loci of $X$ and $Y$, and the linear-section lemma 13.3 occupy sites 76--77. The lemma invokes the Veronese multiples from section 9, while its application uses the Grassmannian incidence notation and Bertini's connectedness statement to construct the elementary fibration (sites 77--78). The proof of the lemma leads to substantial ``remords'' to sections 8 and 9: site 79 cites 2.3 as revised in section 12 and rewrites the linear spaces through a fixed subspace as spaces through a point; sites 79--80 use 8.8 and the symmetry expressed in 8.7(c), then explain the induction on the codimension of a linear section. The handwritten cancellation at site 77 includes the entire clause asserting $C\cap Y=\emptyset$, while the added finite-set avoidance condition remains; at site 78 a whole typed interlinear line belongs after the incidence scheme $\underline{T}^{n-1}_Q$. Remark 13.4 discusses the field hypotheses, records a finite-extension conclusion, treats the omission of geometric irreducibility, and proposes the general-base variant through 7.1 (sites 81--82). Its inserted label c), its renumbering of the following paragraph to d), and the later typed c) and e) show the changing arrangement of these remarks. The finite-set refinement at site 82 and the discussion of the analytic space at site 83 complete the section; the latter describes the homotopy groups, a classifying space $K(\pi_1,1)$, and its universal covering space. Several printed expressions remain especially useful to check beside the scans: definition 13.1 has $X'\to X$ at site 75; the neighborhood $V$ is said to lie ``dans X'' at site 78; site 82 first gives $L_i$ codimension $n-1$ and later writes codimension $n$. The manuscript's own qualifications and unresolved letterforms remain marked, and these expressions have received no mathematical repair.
\par\typeout{EGA-SOURCE-END guide-iv20-13 PDF-PAGE \thepage}

\par\typeout{EGA-SOURCE-BEGIN guide-iv20-14 PDF-PAGE \thepage}
\pdfbookmark[1]{IV 20.14}{source:guide-iv20-14}
\paragraph{IV 20.14: Projections coniques (cote 25, sites 84--91).}
This section supplies the systematic account of conical projection already used at the end of section 8 and in 10.4; its opening note proposes moving the account earlier, perhaps into the auxiliary Grassmannian discussion (site 84). From $C=P(E/G)$ of codimension $m+1$, Grothendieck recalls the Chapter II construction $P-C\to P(G)$ and interprets $P(G)$ as the closed family $Q(C)$ of linear spaces containing $C$. For a morphism $f:X\to P$, he then gives the restricted projection and its extension through the incidence scheme, introducing the notation $\widetilde X(C)$ and $\widetilde F(C)$ and distinguishing the latter $F$ from the quotient used at the start (sites 85--86). The generic-center discussion recalls 5.3 as extended in section 12 and uses 10.2 for the blow-up interpretation. Proposition 14.1, substantially revised at site 87, concerns fiber dimensions; its proof at site 88 returns to section 2 as revised in section 10. Corollaries 14.2 and 14.3 treat finiteness, dominance and the Cohen--Macaulay and flatness conclusions. The added line before 14.3 refers to the remark on manuscript page 21 before 5.6 and writes $\varphi^{(m)}_\xi$, with the small subscript on the following $F$ still uncertain. Site 89 returns to the differential analysis of section 12, preserving the author's qualification ``ou montrerait, si nous ne la faisons pas'', and invokes the section 9 replacement by a Veronese multiple. Proposition 14.4 concerns a finite subscheme $T$: the revised item a) adds the immersion hypothesis and replaces ``un monomorphisme'' by ``radiciel i.e. géométriquement injectif''; item b) receives a further handwritten inverse-image clause. Proposition 14.5 and its added conclusion lead to Corollary 14.6 at site 90, where the original birational-isomorphism wording is replaced by an existential formulation involving a ``plongement birationnel''. The following discussion records limits to the proposed immersion and nonramification conclusions, explicitly mentions Chapter VIII, and explains the varying ramification locus. Site 91 closes with three base-change diagrams and the universal family, returning to 7.1 and proposing $Q$ as the name of a morphism of functors; ``Hilb'' is canceled in its target. The revisions should be read alongside their scans: at site 86 the generic center is literally called a linear subvariety of $X_{k(\eta)}$; in the amended 14.1(b) at site 87 the deletion crosses the printed ``si'' together with the old properness clause; and the dimension condition at site 89 is a typed interlinear insertion whose connector places it after the parenthetical qualification. These page readings and the remaining shorthand doubts are retained without supplying a revised mathematical statement.
\par\typeout{EGA-SOURCE-END guide-iv20-14 PDF-PAGE \thepage}

\par\typeout{EGA-SOURCE-BEGIN guide-iv20-15 PDF-PAGE \thepage}
\pdfbookmark[1]{IV 20.15}{source:guide-iv20-15}
\paragraph{IV 20.15: Axiomatisation de certains des résultats précédents (cote 25, sites 91--95).}
Grothendieck closes IV 20 by considering a common framework for the results of sections 2--8. The opening discussion at site 92 explicitly leaves the extension of Bertini--Zariski unsettled and reports that he has written to Serre and Zariski about it; his suggested differential hypotheses remain part of that unfinished discussion. The framework is a commutative square (D), with a parameter scheme $G$ and an incidence prescheme $\underline P$ over $P$, followed by the corresponding square for $X\to P$. Its assumptions on flatness, geometric irreducibility, relative dimension and the smooth morphism $\underline P\to P$ are set out at sites 92--93. The dimension argument for a closed subset $Z$ leads to a proposed nonemptiness axiom at site 94; the author then states how he expects to recover 2.1--2.11 and formulates a further tangent-space condition for 2.12. The exceptional set $E(x,V)$ is described through smoothness, nonramification and transversality; its designation passes from typed ``constructible'' to an inserted and canceled ``fermé'', then back to handwritten ``constructible''. A Jacobian-criterion and dimension-counting discussion follows, with the quasi-finiteness hypothesis and several closed-set labels added between the lines. The concluding page claims the corresponding variants of 2.12--2.18, reviews the passage from generic to general sections and the uses of sections 5 and 7, and again leaves sections 3 and 6 dependent on the Zariski extension (site 95). Its final instructions to Dieudonné ask for a treatment of section 8 in this framework, while singling out 8.7(c) and its proposed axiomatic reformulation for further work. The guide preserves these as the author's programme and qualifications. In the text, the two signs between $N-n$ and the following $1$ at site 94 remain unresolved plus/overwritten-minus readings; the accompanying arithmetic has not been used to settle the ink. The arrow from the generic-fiber expression $\underline Z_\eta$ to $G$ on the same page, and the overtyped dimension capital after $N=$ at site 95, likewise remain as scan readings with their recorded doubts.
\par\typeout{EGA-SOURCE-END guide-iv20-15 PDF-PAGE \thepage}

\par\typeout{EGA-SOURCE-BEGIN guide-iv21-opening PDF-PAGE \thepage}
\pdfbookmark[1]{IV 21: introduction}{source:guide-iv21-opening}
\paragraph{IV 21, introductory prospectus (cote 25, sites 137--138).}
The separate contents sheet gives four subdivisions: invertible sheaves on a projective bundle and morphisms between projective bundles; representability of the divisor functor associated with an invertible sheaf and relative divisors on projective and multiprojective bundles; linear systems and morphisms to projective bundles; and linear systems with invertible Modules. Its introduction explicitly joins the local work of Chapter IV to the global constructions of Chapter II and some results of Chapter III. Grothendieck presents the language of linear systems as a way to organize both morphisms to projective bundles and invertible Modules on a relative prescheme. He then points toward Brauer--Séveri schemes as parameter schemes, explains their relation to projective bundles after a surjective étale extension, and refers their treatment to descent in Chapter V and to a planned Chapter VI. The canceled phrase tentatively read as ``fibrés principaux homogènes'' before ``torseurs'' records a terminological revision on site 137. Site 138 singles out theorem 1.1 and its first corollaries in subdivisions 1 and 2: its typed replacement describes the central calculation in terms of the Picard group of a projective bundle and that of its base. These are explicit statements of the manuscript's dependencies and intended continuation. See \href{https://grothendieck.umontpellier.fr/web/viewer.html?file=../25.pdf#page=138}{site 137} and \href{https://grothendieck.umontpellier.fr/web/viewer.html?file=../25.pdf#page=139}{site 138}.
\par\typeout{EGA-SOURCE-END guide-iv21-opening PDF-PAGE \thepage}

\par\typeout{EGA-SOURCE-BEGIN guide-iv21-1 PDF-PAGE \thepage}
\pdfbookmark[1]{IV 21.1}{source:guide-iv21-1}
\paragraph{IV 21, \S 1: invertible sheaves on projective bundles and their automorphisms (cote 25, sites 98--105).}
The section begins with the classification of invertible Modules on $P(\underline E)$ and develops it into a description of morphisms and automorphisms of projective bundles. Theorem 1.1 and Corollaries 1.3--1.5 describe the partition of the base by the integer twist and the remaining invertible Module on the base; Remark 1.2 separates the rank-zero, rank-one and higher-rank parts. The proof first establishes uniqueness, then reduces existence to the field case and invokes III 4.6.5 with the vanishing from III 2; its preliminary constructions and the canonical projective-bundle isomorphisms refer to Chapter II. At site 101 the author explicitly revises the exposition of the field case, asks to begin directly with the graded module of sections, and replaces the connected-base version by the formulation using locally constant integer-valued functions $\underline Z(S)$. The resulting map from the constant group scheme to the Picard scheme is assigned to the planned Chapter V, at sites 101--102. Site 102 makes a further revision: the author proposes stating 1.1 without its initial rank restriction, separates the uniqueness assertions according to nonempty fibers and positive-dimensional fibers, and moves Remark 1.2 into the proof. The same page introduces degree, and sites 102--104 describe morphisms of a fixed degree through an invertible Module $\underline M$ and a homomorphism involving $\underline{\mathrm{Sym}}^n(\underline E)$. The cases of degree 0 and degree 1 are identified in the manuscript as the constant and the linear morphisms. Theorem 1.6 and Corollaries 1.7--1.8 treat isomorphisms and their local descriptions; Remark 1.9 names the projective group and its absolute version $GP(r)$. Its handwritten instruction at site 105 asks that the representability argument be carried out using a determinant open; the position of the possible dual accent in that marginal formula is above the first $E$, and its presence remains doubtful. Several revisions require the reader's attention without changing the text: site 98 prints $\underline O_P(1)$ next to a final twist $(n)$ in the formula of 1.1; site 101 retains the printed $q$-operators in the proposed field-case exposition; and site 102 says ``$\underline L$ inversible sur $S$'' in a reformulation whose preceding discussion concerns invertible Modules on $P$. The overprinted parenthesis/prime-like mark in the inverse image formula at site 103 remains uncertain. Large diagonal marks occur at sites 98--103; their graphic form is recorded, but their function does not securely establish deletion of entire arguments. In particular, the unsupported paragraph-wide deletion wrappers at site 100 have been removed, and the insertion connector at site 101 is not treated as canceling ``comme Module''. The page before this section, site 97, contains a severely faint related text whose uncertain readings remain marked after inspection; it has supplied no reading for these clearer pages.
\par\typeout{EGA-SOURCE-END guide-iv21-1 PDF-PAGE \thepage}

\par\typeout{EGA-SOURCE-BEGIN guide-iv21-2 PDF-PAGE \thepage}
\pdfbookmark[1]{IV 21.2}{source:guide-iv21-2}
\paragraph{IV 21, \S 2: relative divisors on projective and multiprojective bundles (cote 25, sites 106--111).}
The section applies the classification of invertible Modules from 1.1 to positive relative divisors. At sites 106--107 the canonical direct-image formula marked $(*)$ converts the section defining a divisor into a map involving a symmetric power of $\underline E$; the author expresses fiberwise regularity through universal injectivity and the surjectivity of its transpose. Proposition 2.2 describes the resulting correspondence with invertible quotients of the dual symmetric power, and Corollaries 2.3--2.4 give the representing projective bundles and their sum over the degree. Paragraph 2.5 introduces a finite family of projective bundles and their fiber product. Proposition 2.6 and paragraph 2.7 extend the classification and degree notation to this multiprojective bundle, explicitly treating the partial degrees as locally constant functions on the base. Sites 108--109 connect this construction to the planned paragraph 25 on the resultant and discriminant of forms. Paragraph 2.8 invokes Chapter II, paragraph 2, in the displayed direct-image formula $(xx)$; its canceled alternative reference to Künneth is readable. Proposition 2.9 and Corollary 2.10 describe divisors of a fixed multidegree by invertible quotients of a tensor product of dual symmetric powers and give the associated representing projective bundle and sum. Paragraph 2.11 then extracts the relative representability argument from the projective-bundle case and relates it to the Picard prescheme. At site 110, under the displayed hypothesis on $f_*(\underline O_X)$, the proof uses Chapter III, paragraph 7, to introduce the finitely presented Module $\underline Q$ representing the relevant functor, identifies sections with homomorphisms to an invertible Module, and states the correspondence with sections of $\underline P(\underline Q)$. The proposition number on that page is typed 2.11 with an uncertain pen correction resembling 2 on its last digit; site 111 explicitly refers to 2.12. Theorem 2.13 at site 111 gives the open retrocompact representing subprescheme, and the text identifies the entire projective bundle in its geometrically integral-fiber case. Its proof invokes the openness of transverse regularity through the incomplete reference ``par.11 ...'' and obtains retrocompactness by reduction to the noetherian case. The transcription retains several visible features that affect reading: site 106 says ``section de $\underline E$'' and uses $\underline O_X(\underline D)$; site 107 prints ``décomposition de D'' and has a crowded extra subscript and crossed group in the symmetric-power expression of 2.4; the displayed isomorphism of 2.6 at site 108 omits the $\underline M$ factor mentioned in the surrounding statement; and site 109 prints $\underline O_X(D)$ and $f_*(\underline D)$ in the two formulas. These are preserved with the scanned alphabets and indices. The N/G variable at site 110, the corrected proposition digit, and several typewriter overprints remain explicitly uncertain.
\par\typeout{EGA-SOURCE-END guide-iv21-2 PDF-PAGE \thepage}

\par\typeout{EGA-SOURCE-BEGIN guide-iv21-3 PDF-PAGE \thepage}
\pdfbookmark[1]{IV 21.3}{source:guide-iv21-3}
\paragraph{IV 21, \S 3: linear systems, fixed points and maps to projective bundles (cote 25, sites 112--122).}
The handwritten heading at site 112 introduces linear systems of divisors and morphisms into projective bundles; its large section digit retains overlapping 3/2 traces. The opening preliminary account of relative divisors and their representing functor is visibly canceled. The following live text defines fixed points of a family and studies the open complement of the fixed-point set. Sites 113--114 compare the family on $X/S$ parametrized by $T$ with the transposed family on $T/S$, including the hypotheses on flatness and geometrically irreducible fibers, and then discuss fixed points ``au sens large''. The author separates an added sorite about the transposed family, marks its relocation to paragraph 19, and adds a qualification about its behavior under base change. At site 115 the projective-bundle representation of the divisor functor supplies the degree of a family and the definition of a linear system. The associated morphism on the complement of the fixed points is related to the canonical incidence divisor $\underline H$ on $P\times_S\check P$. Proposition 3.1 at site 116 states the correspondence for systems without fixed points and expresses its condition through the images of associated prime cycles after extension to an algebraic closure. The handwritten additions to the preceding discussion include the non-degeneration condition and the terminology of a pseudo-morphism; the text refers back to paragraph 11 and proposes moving part of the discussion to paragraph 19.

Site 117 defines fixed components and gives a uniqueness proposition whose original typed number 2 has an unreadable handwritten replacement. Its proof compares two systems with the same pseudo-morphism, uses schematic density and replaces the earlier SGA reference with \S 20.10; a marginal reference points to \S 20.11. Lemma 3.2 at site 118 gives the direct-image extension statement under the printed flatness and depth conditions. Its proof reduces to the noetherian case and cites paragraphs 6 and 5, with an authorial suggestion that the lemma should appear in paragraph 11. The page then completes the preceding uniqueness argument using the maximal domain of the pseudo-morphism. The long bracket and wavy pen line spanning the opening argument and lemma have unresolved editorial function, so their body text is retained. Two conspicuous printed features also remain: the quantifier in the lemma reads $x\in X$, and the displayed immersion has target $T$ in $i:X-T\to T$.

The revised proposition at site 119 returns to this correspondence using two sets, written with distinct oval-stemmed and forked-stemmed handwritten glyphs. Their alphabetic identification remains uncertain; the transcription preserves that distinction. Its heading has adjacent 3-like figures and an uncertain underlying 2, without a securely written separator. The statement relates systems satisfying the displayed depth condition to pseudo-morphisms and gives an invertibility criterion involving $i_*(\underline L_U)$, with $\underline L_U$ obtained from the inverse image of $\underline O_P(1)$. The inserted non-degeneration condition refers explicitly to 3.1. Sites 120--121 continue the proof by flat base change to the dual projective bundle, tensoring with an invertible Module and comparison with the incidence divisor. The argument cites II 7 and the extension lemma. Site 120's main-body formula for $D'$ lacks an inverse-image exponent, while its marginal note writes that exponent explicitly; both written forms are preserved. The long marginal note remains partly cut and partly illegible. At site 121 the first inserted lemma reference reads 4.2, whereas the later references on the same page read 3.2. Corollary 3.4 states the restriction and extension assertions for locally free Modules and divisors under the stated depth hypothesis, including uniqueness and schematic closure in the added note. Its clause (b) visibly says that $\underline L'$ is on $X$. Corollary 3.5 and its continuation at site 122 use local factoriality, with regularity as the example, to obtain the asserted bijection; the final handwritten qualification excludes factorization through a hyperplane after passage to $\bar k$. The proof invokes Auslander--Buchsbaum and replaces its planned Chapter III reference with \S 20. The older final paragraph at site 122 discusses an invertible Module associated with a linear system. It lies beneath crossing diagonals with an arrow-like terminal whose function is unresolved, so its words and formulas remain available. Its blank relation positions, typed superscript $x$, and uncertain decoration on the last $P$ are retained without completing the formulas.
\par\typeout{EGA-SOURCE-END guide-iv21-3 PDF-PAGE \thepage}

\par\typeout{EGA-SOURCE-BEGIN guide-iv21-4 PDF-PAGE \thepage}
\pdfbookmark[1]{IV 21.4: linear systems}{source:guide-iv21-4}
\paragraph{IV 21, subdivision 4: linear systems of divisors and invertible Modules (cote 25, sites 123--133).}
The heading on site 123 replaces a typed 7 by a handwritten 4. The opening calculation explicitly invokes subdivision 1, replaces a reference to subdivision 2 by \S20.3, and later cites Chapter III, \S2. It describes a divisor on the product with the dual projective bundle by an invertible Module \(N\) and a section \(\phi\), with fiberwise regularity and a degree-one requirement. The calculation expresses \(N\) through an invertible Module \(L\) on \(X\), identifies \(L\) by a direct image of \(N(-1)\), and uses projection formulas to pass to a homomorphism \(g^*E\longrightarrow L\), or equivalently \(u:E\longrightarrow g_*L\). Several overwritten projection indices remain uncertain in the transcription. Proposition 4.1 on site 124 organizes this construction as a correspondence with pairs \((L,u)\); its regularity test is stated after field extension, and the geometrically integral-fiber case is expressed by injectivity of the corresponding map into \(H^0(X_s,L_s)\). A marginal addition requires nonempty projective fibers. The author also requests an explicit reminder of the divisor obtained from the section of \(L\otimes\mathcal O_{\check P}(1)\). Corollary 4.2 begins on the same page with a finitely presented representing Module \(Q\). Its continuation on site 125 cites III 7.7.6 for the existence of \(Q\), uses Nakayama to pass from fiberwise surjectivity to a quotient \(Q\longrightarrow\check E\), and describes the remaining equivalence by multiplication with invertible sections. An added note records the role of \(f_*\mathcal O_X\simeq\mathcal O_S\) and of the corresponding units. The alternative description uses the divisor functor of subdivision 2 and the morphism \(\check P=P(\check E)\longrightarrow P(Q)\); site 126 then presents the associated classes through the Grassmannian of \(Q\). The manuscript's internal reference ``comme dans 5.1.'' in 4.2, the printed ``sur X'' for its earlier projective bundle, and its changes between \(f\) and \(g\) are retained in the text. See \href{https://grothendieck.umontpellier.fr/web/viewer.html?file=../25.pdf#page=124}{site 123}, \href{https://grothendieck.umontpellier.fr/web/viewer.html?file=../25.pdf#page=125}{site 124}, and \href{https://grothendieck.umontpellier.fr/web/viewer.html?file=../25.pdf#page=126}{sites 125--126}.

Proposition 4.3 begins on site 126 and treats the morphism from the parameter bundle to \(\mathrm{Div}_{X/S}\), followed by the action of automorphisms of that parameter bundle. The first part uses an associated prime cycle and its global functions; the second part is formulated with surjectivity of the structural map and includes a handwritten functorial addition. The author supplies a concrete warning on site 127: for an open subscheme of \(P^1_k\) omitting two distinct rational points, those two points give the same zero divisor in the displayed family. The proof first works over an algebraically closed field, passes to an associated cycle with the reduced induced structure, and compares sections differing by invertible functions. The lemma on site 128 chooses \(r+1\) points whose images form a projective basis. Its use on site 129 is explicitly tied to subdivision 1's description of automorphisms of a projective bundle. The general-base argument then passes through the noetherian case and local artinian quotients. On sites 130--131 the difference of two sections is expressed through tangent vectors and differentials, with separate descriptions for the divisor functor and for the morphism functor. On site 130, the written definition \(F_o=F\otimes_A k\) is followed by the phrase ``fibre réduite''; the argument also says \(h_o\) is surjective immediately after naming the preceding homomorphism \(h_o^*\). Site 131 gives the Hom expression the subscript \(\underline O_{X_o}\) and the target \(\underline O_{P_o}\). These literal choices remain in the text. Grothendieck says that this part should refer to the calculations of the planned paragraph on infinitesimal extensions, whose number is being changed from 25 toward 26. He also records the dependence of those calculations on paragraph 16 and explains the intended order of the references. Sites 131--132 describe the resulting interpretation of a linear system as a subfunctor of the divisor functor, the uniqueness of the parameter-bundle isomorphism, and the proposed use of faithfully flat descent in Chapter V with twisted projective bundles. These are the manuscript's own proof steps and plans; the guide supplies no missing tangent-space argument. See \href{https://grothendieck.umontpellier.fr/web/viewer.html?file=../25.pdf#page=128}{sites 127--129}, \href{https://grothendieck.umontpellier.fr/web/viewer.html?file=../25.pdf#page=131}{sites 130--131}, and \href{https://grothendieck.umontpellier.fr/web/viewer.html?file=../25.pdf#page=133}{site 132}.

The closing proposition on sites 132--133 is referred to in the surrounding discussion as 4.4, although its displayed number has been canceled. It returns to the set \(\mathrm{Ens}(D)\) of divisors obtained from rational points of the parameter bundle. The page replaces a description by rational functions with one involving a vector subspace of meromorphic functions and regularity of its nonzero members; the successive cancellations and the later occurrence of ``pseudo-'' remain visible. The last part compares the subspaces by multiplication with a regular function and records the remaining scalar ambiguity. Grothendieck leaves the proof as an exercise using 4.1 and proposes moving this proposition ahead of 4.3. He then discusses the additional algebraic-closure condition and its geometric-integrality case, before adding the convention that an ensemble of the form \(\mathrm{Ens}(D)\) may also be called a linear system of divisors. The nearby marginal discussion of complete linear systems on site 126 remains partly clipped and uncertain; it is not completed from this closing proposition. See \href{https://grothendieck.umontpellier.fr/web/viewer.html?file=../25.pdf#page=133}{site 132} and \href{https://grothendieck.umontpellier.fr/web/viewer.html?file=../25.pdf#page=134}{site 133}.
\par\typeout{EGA-SOURCE-END guide-iv21-4 PDF-PAGE \thepage}

\par\typeout{EGA-SOURCE-BEGIN guide-letter PDF-PAGE \thepage}
\pdfbookmark[1]{Letter to Dieudonné}{source:guide-letter}
\paragraph{Letter to Dieudonné, 29 September 1965 (cote 25, sites 134--135).}
The handwritten date fixes this letter to 29 September 1965. Grothendieck thanks Dieudonné for a contents list for paragraphs 16--19 and asks for the provisional contents of 20 and 21, accepting their inclusion in the fourth fascicle of Chapter IV while asking how his former paragraph 20 will be divided. The proposed continuation assigns linear systems and additions on the Picard group to 22, Grassmannians to 23, smooth forms and ordinary quadratic singularities to 24, hyperplane sections to 25, the resultant and discriminant to 26, and infinitesimal extensions to 27; 20 and 21 are still represented by question marks. His joke about the pleasing number $27=3^3$ accompanies an expectation that the long paragraph on hyperplane sections will be divided again. These are the numbers of the plan recorded by this letter. The scan supplies no identification of them with every later number used in the surrounding prenotes. The second subject is authorship: he declines to publish the former appendix to paragraph 18 under his name alone, credits Dieudonné's extensive rewriting, and asks him to reconsider a joint paper. Site 135 also preserves a completely canceled paragraph referring to ``par.20, 10.9.1.'' and to paragraph 12: it invokes constructibility, and then openness under the stated flatness and finite-presentation assumptions, of a fiberwise depth locus, as part of the familiar argument for passing farther along the indexing system. The handwritten deletion belongs to the letter's textual history and remains visible in its diplomatic form. The check confirms the date, the cancellation, and the proposed contents directly from the two scans; it establishes no bibliographic identification with the separately named L49d. See \href{https://grothendieck.umontpellier.fr/web/viewer.html?file=../25.pdf#page=135}{cote 25, site 134} and \href{https://grothendieck.umontpellier.fr/web/viewer.html?file=../25.pdf#page=136}{site 135}.
\par\typeout{EGA-SOURCE-END guide-letter PDF-PAGE \thepage}

\par\typeout{EGA-SOURCE-BEGIN guide-planning PDF-PAGE \thepage}
\pdfbookmark[1]{Planning sheet and rotated fragment}{source:guide-planning}
\paragraph{A further contents note and a rotated fragment (cote 25, site 136).}
This handwritten sheet places additions on rational maps, birational morphisms and blow-ups under paragraph 20; Cartier and Weil divisors, Picard, linear systems and Brauer--Severi under 21; Grassmannians under 22; hyperplane sections under 23; and infinitesimal extensions under 24. It therefore supplies another explicitly numbered plan alongside the dated letter. Its lower fragment is written with the sheet turned through 180 degrees. The visual check recovers the labels ``poids $n$,'' ``poids $p+q$,'' and ``filtration $p$,'' with the last $p$ doubly underlined, around a cohomological formula; the exponent of the first cohomology group is a superscript $n$, while another overwritten exponent remains doubtful. A crossed heading asks for a ``signification'' and ends ``intrinsèque,'' with its intervening abbreviation still marked uncertain. The scan gives neither a date nor an explicit connective statement tying this fragment to the contents above, so the edition preserves the fragment and its orientation without supplying such a statement. See \href{https://grothendieck.umontpellier.fr/web/viewer.html?file=../25.pdf#page=137}{cote 25, site 136}.
\par\typeout{EGA-SOURCE-END guide-planning PDF-PAGE \thepage}

\par\typeout{EGA-SOURCE-BEGIN guide-blass PDF-PAGE \thepage}
\pdfbookmark[1]{Blass's translation and its extent}{source:guide-blass}
% Completed guide assembled as the successive parts of the cote24 audit were finished.
\subsection*{Piotr Blass's translation: extent, opening, and notation}
The 1987 typescript consists of a title page, Summary, contents, and the English text of IV 20, Sections 1--11, on archive sites 4--74. Its contents list further sections, but the translated mathematical body ends with the Grassmannian discussion corresponding to French site 65. The handwritten French note on site 75 is retained in all three Blass readers; the stray CV sheet on site 76 is retained only in the diplomatic reader. These books preserve Blass's English, his French-language glosses, his typed questions, and the visible handwritten annotations. The following comparisons come from the two checked witnesses, and the decision records locate each gap or query by its French archive site. Sources: \href{https://grothendieck.umontpellier.fr/web/viewer.html?file=../24.pdf\#page=2}{cote 24, sites 1--76}; \href{https://grothendieck.umontpellier.fr/web/viewer.html?file=../25.pdf\#page=3}{cote 25, IV 20}.

The English title page identifies the work as ``Stage 0 -- translation by Piotr Blass,'' and the following Summary places the translated \emph{Commentaires} before the contents. Its queries belong to this witness: ``reorganize???'' corresponds to \emph{réorganisation de fond en comble} on French site 2; the typed question about no.\,16 and its handwritten ``Ask AG.'' refer to the French statements on sites 2--3. In Section 1, the English develops the same quotient-Module, hyperplane, relative-divisor, and incidence constructions as French sites 4--8. The opaque strips on English sites 4 and 6 remain unreadable, with their positions identified rather than completed from the French. English site 5 asks which paragraph 17 is meant and queries the separate circle above the functor domain; both corresponding passages are on French site 5. English site 7 retains four unfilled arrow positions and asks whether \emph{Module} always means a coherent or quasi-coherent sheaf of modules, at the passage corresponding to French site 7. Its drawn diagram contains four connectors without arrowheads. The comparison also records exact differences of notation: the final term of the transpose formula on English site 4 is $\underline O_P(1)$, while French site 4 types $\underline D_P(1)$ and then names a section of $\underline L_P(1)$; English site 4 types an ordinary $p*$ where the French has a handwritten subscript star; English site 5 lacks the French dual marks in the symmetric-power expression and on the final projective space; English site 6 lacks one dual mark in the section over $S'$; and English site 7 types the continuous string $XPV$ where French site 7 says that $F$ is a Module on $X$. These visible forms are preserved. Full-height baseline letters and genuinely raised characters have been compared separately, since the typescript uses both. Sources: \href{https://grothendieck.umontpellier.fr/web/viewer.html?file=../24.pdf\#page=2}{cote 24, sites 1--8}; the corresponding French passages are \href{https://grothendieck.umontpellier.fr/web/viewer.html?file=../25.pdf\#page=3}{cote 25, sites 2--8}.

\paragraph{Blass's Section 2 (sites 9--19 opening).}
The translation follows the local study from irreducibility and finiteness through associated cycles, regularity, smoothness, and separability, and retains the author's comments on reorganizing the results. Its queries can be located precisely: site 9 questions the insertion \emph{En tous cas} in French 8; site 10 asks about underlined capitals; site 12 queries the added alternative in Proposition 2.6 and the paragraph-3 reference of French 9--10; site 13 leaves the marginal remark about the \emph{dim.\,|\,dém.} of 2.1 unread and retains four reference placeholders from French 11; site 15 leaves the argument of $k(\ )$ blank, as French 13 does; and site 17 asks about the two paragraph-6 references on French 13--14. Both opaque masks on English 16 remain unreadable. The first ``PB incomprehensible'' on English 18 corresponds to the repeated $Z_i$ in French 14's first assertion of 2.16; the repetition is present in both witnesses. Its second query accompanies the reduction to an algebraically closed field, whose order relative to the separability sentence differs from the French insertion. The mathematical comparisons retain all observed differences: English 9 inserts $k$ in $\operatorname{Spec}k(\eta)$, where French 8 has $\operatorname{Spec}(\eta)$; English 10 types $Z_\eta$ in the proof where French 8 corrects the letter to $Y$; English 11 types $Z1$ where French 9 has $Z'$; English 13 types a capital $N$ in the coprofundity bound and $H_\eta$ in the conclusion, where French 11 has lowercase $n$ and the doubtful $E|F$ with indices $\eta,y$. English 14 places its first bar over \emph{dim}, omits the inverse-image exponent in $h(Z)$, and says that constructibility is already known by applying 2.3, while French 11--12 bars $f(Z)$, writes the inverse-image exponent, and says that constructibility must be known to apply 2.3. English 15 has a lowercase $y$ and ordinary arrow into $U_{k(\eta)}$ against French 12's capital $Y$ and hooked arrow, and retains the baseline $1$ and lowercase $u/k$ in its differential notation. English 16 places the bound $1\leq i\leq n$ beneath the quotient sum; French 13 places it after the inserted generators $a_i$. English 16--17 omit the French marginal qualification beside 2.14, whose first term remains doubtful, while English 17 types $k(y)$ in the equivalence and omits ``geometrically'' before the first ``normal'' in 2.15. Capital $K(\eta)$ occurs on English 16--17 where the corresponding French uses lowercase $k$. Several additional underlines on $Y$ are visible on English 9, 11, 12, and 14. English 18 replaces the French ranges 2.13--2.17 and 2.12--2.17 by conjunctions of their endpoints; the explanation continues on English 19, with an extra typed arrow before ``smooth = regular.'' None of these forms is repaired from the other witness. Sources: \href{https://grothendieck.umontpellier.fr/web/viewer.html?file=../24.pdf\#page=10}{cote 24, sites 9--19}; \href{https://grothendieck.umontpellier.fr/web/viewer.html?file=../25.pdf\#page=9}{cote 25, sites 8--15}.

\paragraph{Blass's Section 3 (sites 19--21).}
The translation presents the Bertini--Zariski theorem through the function-field description and Zariski's lemma, retains the request that Dieudonné choose a proof from the literature, and follows with the geometric-integrality and connectedness corollaries. The numbering query on site 19 corresponds to the overwritten heading on French 15; ``[PB: check this!]'' points to the inclusion described by sending $S_i$ to $S_i$ on French 16. Several literal formulas require care. English 19 has $\dim\overline{f(X)}\geq0$ where French 16 has $>0$, and ends $L=K(S_1,\ldots,S_n)$ at $n$ where the French ends at $n-1$. English 20 uses $a_iS_1$ where French 16 has $a_1S_1$, and omits the minus before $t_it_n^{-1}$ in the definition of $a_i$. Its lemma omits the French definition $L=K(S_1,\ldots,S_m)$, while the sum has printed capital limits $I$ and $M$ against French 17's $1$ and $m$; the English field name is the literal baseline $k1$. English 21 uses a clear typed $\sigma$ and $f(X)$ where the checked French 17 retains doubts about the script family letter and the argument $Z|X$. It also prints capital $K(\eta)$ in the conclusion and ``$X_0$ and $X'$'' at the start of the chain, against the French lowercase field letter and equality. Its handwritten query points to the paragraph-6 reference accompanying Hartshorne's theorem. The guide reports these readings and retains the French uncertainties. Sources: \href{https://grothendieck.umontpellier.fr/web/viewer.html?file=../24.pdf\#page=20}{cote 24, sites 19--21}; \href{https://grothendieck.umontpellier.fr/web/viewer.html?file=../25.pdf\#page=16}{cote 25, sites 15--18}.

\paragraph{Blass's Section 4 (sites 21--25).}
The translation returns to the general base and uses constructible properties to pass from the generic hyperplane to sufficiently general hyperplanes, then lists the Module and prescheme properties of 4.2--4.4. The doubtful handwritten query on English 22 points to the paragraph-9 reference in French 18; English 25 asks both what AQT means and where paragraph 9 is, at the opening of French 21. English 22 prints $G$ at a hypothesis where French 18 retains the letter doubt $F|E$, omits ``two'' from ``the following sections,'' and inserts a further $\xi$ before ``$k$ algebraically closed.'' English 23 uses ``scheme,'' \emph{coprof}, and ``Cohen Macauley'' where French 19 has \emph{préschéma}, \emph{coppof}, and ``Cohen-Macaulay''; its characteristic character is the round letter O, distinguished on the same page from the taller numeral 0. The wording ``F is satisfied at P'' is also literal. English 24 lists geometric integrality before geometric irreducibility, following the direction of the French relocation while the overwritten French labels remain uncertain. More consequentially, its characteristic-zero and unramified clauses are joined by ``and,'' against French 20's \emph{ou}; the French inserted morphism description itself still carries a reading doubt. English 25 prints $K(\xi)$ against lowercase $k(\xi)$, and $f(X)$ against $f(X_i)$. Its $E_\xi$ and final projective-space symbol without a dual are already present in French 20. The comment following 4.4 says ``not valid as such if f is quasi-finite,'' whereas French 21 says \emph{n'est valable tel quel que si f est quasi-fini}; the next instruction becomes ``it is sufficient to consider'' where the French says \emph{il faut regarder}. The comparison records these exact changes in qualification without repairing the statements. The Latin instruction is retained with Blass's added label ``[Latin].'' Sources: \href{https://grothendieck.umontpellier.fr/web/viewer.html?file=../24.pdf\#page=22}{cote 24, sites 21--25}; \href{https://grothendieck.umontpellier.fr/web/viewer.html?file=../25.pdf\#page=19}{cote 25, sites 18--21}.

\paragraph{Blass's Section 5 (sites 25--29).}
The opening incorporates the paragraph on the openness of $U$ and $E$ that French 21 cancels and French 22 explicitly restores as text struck out in error. The query ``PB ???'' and the handwritten ``illegible'' on English 25 refer to its definition of $\xi$ as the image of $y$; English 26 then completes the paragraph and develops the section defining the universal hyperplane divisor. The following regularity and openness statements correspond to French 22--24. English 26 leaves a star after ``Mike'' without a footnote and question punctuation in the parenthesis about $S'$ over $S$; French 22 has \emph{comme dirait Mike} and canceled doubtful question marks in that parenthesis. English 27 leaves a question at $\mathrm{id}_F\otimes$ where French 22 has $\mathrm{id}_F\otimes\varphi$, and a further star at the paragraph-3 reference of French 23. Its incomplete $o_{\mathrm{III}}$ reference corresponds to French $E_{\mathrm{III}}$. In the proof of 5.3(b), English 27 prints equality between the intersection and the repeated $Z$ term, against French 23's handwritten inequality; its final section symbol also lacks the French $\eta$ index. English 28 has ``P is flat over S'' against the inserted French ``F,'' and ``a little less evidently'' against \emph{évidemment}. It asks for the references 12.1.1 and 12.1.4 cited at the start of French 24's 5.6. English 29 leaves blank the relation between $U_P$ and $U$, where French 24 inserts inclusion, omits $\xi$ after the parenthetical $Y$, and lacks $\varphi_\xi$ in the regularity condition of 5.7. The incomplete bracket ``[fascicule'' also remains. Sources: \href{https://grothendieck.umontpellier.fr/web/viewer.html?file=../24.pdf\#page=26}{cote 24, sites 25--29}; \href{https://grothendieck.umontpellier.fr/web/viewer.html?file=../25.pdf\#page=22}{cote 25, sites 21--24}.

\paragraph{Blass's Section 6 (sites 30--31 opening).}
The translation combines the generic connectedness statement 3.3 with the cited 15.6.3 in Proposition 6.1 on arbitrary hyperplane sections, and retains the planning remarks about 4.4, sections of excessive dimension, and grouping the connectedness questions in one number. The query ``[x X proper over k suppose]*'' on English 30 corresponds to the insertion ``X propre sur S'' in French 25; there is no starred footnote on that English page. English 30 leaves $f(Z)$ unbarred where French 26 has a handwritten closure bar, and prints ``if Y is proper over k'' where the French gives \emph{car Y est propre sur k}. Its three long diagram arrows are drawn in pen and are marked accordingly in the checked transcription. English 31 adds ``(or section)'' after ``No.'' in the planning sentence. Sources: \href{https://grothendieck.umontpellier.fr/web/viewer.html?file=../24.pdf\#page=31}{cote 24, sites 30--31}; \href{https://grothendieck.umontpellier.fr/web/viewer.html?file=../25.pdf\#page=26}{cote 25, sites 25--26}.

\paragraph{Blass's Section 7 (sites 31--35).}
This section follows the use of rational points over infinite fields and étale base changes over finite residue fields to construct hyperplane sections and multisections, then the statements 7.2--7.4 and Remark 7.5. It retains the references to 4.1--4.3 and no.\,5 and the induction on fiber dimension in 7.4. English 31 types the literal section sign \S\ for the section variable where French 27 has $\xi$, and initially has $\xi$ without the French inserted index $0$, although $\xi_0$ appears later in the English paragraph. English 32 asks about ``closer and closer (by successive approximations ???)'' for French 28's \emph{de proche en proche}, with a handwritten ``Translator's note.'' Its printed $k1$ lacks the denominator of the French residual extension $k'/k$. English 33 repeats ``finite'' in the presentation-and-flatness hypothesis and uses $\xi,s$ as the indices of $Y$ and $F$, against the French $\xi,\xi(s)$. The neighborhood in $E$ occurs in both witnesses; the French range in the reference to 4.3 remains doubtful. English 34 leaves several placeholders at French 29's added irreducible-component condition and the references 12.2.1(x) and (xi); part of the French added N.B. also remains uncertain. ``Text crossed out'' represents the canceled paragraph continuing into French 30's Stein-factorization discussion. The underlined ``X [blank] U'' query in 7.4 remains unresolved by French 30, which likewise shows no intervening sign. The raised $1$ in a later $S^1$ is retained alongside baseline $S1$ occurrences. English 35 omits ``over'' before the final $S$ in the nonprojective case and leaves ``ILLEGIBLE'' for the French marginal suggestion about avoiding an \emph{ensemble} whose adjective remains doubtful between \emph{fini} and \emph{donné}. Sources: \href{https://grothendieck.umontpellier.fr/web/viewer.html?file=../24.pdf\#page=32}{cote 24, sites 31--35}; \href{https://grothendieck.umontpellier.fr/web/viewer.html?file=../25.pdf\#page=28}{cote 25, sites 27--30}.

\paragraph{Blass's Section 8: exceptional hyperplanes (sites 36--42).}
The section begins with the dimensions of exceptional hyperplane loci, then the coprofundity statements and the regularity criteria 8.4--8.6. English 36 repeats the clause defining the set of $\xi$ and says ``a better upper bound'' where French 31 changes \emph{majoration} to \emph{minoration}; the query asking for a number corresponds to the French's own unfinished \emph{cf n°}. English 37 asks whether a numerical hypothesis is two, at French 32's $\dim f(X_i)>0$, and uses $G_{k(\xi)}$ in its first coprofundity expression where French 32 has $F_{k(\xi)}$. It omits the intervening assertion about $\operatorname{coprof}G_\xi$ and prints an incomplete $O_{\mathrm{IV}}$ reference against French $E_{\mathrm{IV}}$. Its later gap corresponds to French 33's assertion about $G_\xi$ satisfying $(S_{k-1})$. The extensive queries on English 38 point to French 33's partly unreadable opening, its clear $k=2$, the added reference to 8.5, and the unfinished handwritten discussion at the foot. English 38 says the $(S_2)$ conditions are met where the French says \emph{ne sont pas}. English 39's question about an insertion follows the component clause developed in French 34's relocated continuation, with its zero-dimensional image and quasi-finite case; its lemma retains the literal bound $\operatorname{codim}(Y,X)\geq K$. English 40 leaves blank the relation that French 35 gives as $\operatorname{codim}(T_i,X_j)=k$ and omits the intervening alternative involving $\dim f(X_j)>0$, a point outside $Z_0$, and the next exceptional clause. Its local-ring indices $Xj,Ti$ share one lowered tier, and the production layer of the inline slashed C remains uncertain. English 41's gaps and ``floppy or perhaps sloppy'' comment accompany an already incomplete French 35 statement; its later ``R of'' gap corresponds to French 36's $X$, and its ``finite'' differs from French \emph{quasi-fini}. English 42 omits the dimension parenthesis, \emph{est vide}, and a displayed coprofundity bound from French 36, then prints ``by preference'' for \emph{à priori}, $Z=X$ for $z=x$, and a depth bound $\geq k$ where the French has $<k$. Its final editorial sentence changes the instruction about stating a corollary with 8.6, and the new differential discussion begins with $f:X\to f$ against French 37's $f:X\to P$. The footer query has no matching body star, so its exact scope remains unsettled. Sources: \href{https://grothendieck.umontpellier.fr/web/viewer.html?file=../24.pdf\#page=37}{cote 24, sites 36--42}; \href{https://grothendieck.umontpellier.fr/web/viewer.html?file=../25.pdf\#page=32}{cote 25, sites 31--37}.

\paragraph{Blass's Section 8: conormal calculations and 8.7--8.9 (sites 43--50).}
The conormal Module, universal singular-zero locus, and incidence interpretation lead into the differential diagram and Theorem 8.7. The query about the last two symbols on English 43 points to the kernel of $E_X\to O_X(1)$ on French 37; the queries on English 44 point to French 38's clipped \emph{comm}, \emph{ledit sous-espace vectoriel}, and linear subvarieties containing $a$. The conormal letter retains its unresolved alphabet, and the English section symbol after ``image by'' lacks the French prime. English 45 prints $dt(r-d-1)$ where French 39 has $d$ followed by a doubtful plus, asks about paragraph 16, and leaves a query after the differential-section formula $\psi=d\phi|Y$. Both English $O$ indices are lowercase $y$, the second underlined; both $P$ letters in the first $\Omega$ index occupy a deeper lowered tier. English 46's diagram has typed objects and pen connectors, including one connector without a secure head and two dashed diagonals. Its $B$ and partial-shaped labels differ from French 40's $\beta$ and $\alpha$; its top object, the copied uncertain pen letter, and the missing lower-left base letter remain unresolved. The question about $\psi$ corresponds to French 40's inserted $\mu$. Its displayed composition goes directly to the $(1,1)$ target, whereas the French first writes $(1,0)$ and then interprets the composition as a section with $(1,1)$. English 47 retains the duplicate ``above lies over'' and joins the announced result to the next editorial paragraph. English 48 preserves the paragraph-16 and definition-reference queries and leaves the first projection unread, where French 41 gives $x$; its $H$ bears a raised numeral 1 against the French prime. More substantially, English 48 jumps from ``image of the tangent map of'' to ``the closure of its image,'' omitting the end of 8.7(c), its displayed map, the proof and diagram (D'), and the heading and opening hypotheses of Corollary 8.8 on French 42. The later conjunction ``(i) and (ii) (bis)'' replaces the French range. English 49's queries correspond to French 43's N.B. about the immersion hypothesis and the inserted \emph{ou}; it prints $\leq r-2$ where the French changes the sign to equality, $T$ for $Y$ in the image clause, twice ``(ii ter)'' for ``(ii bis),'' unbarred capital $K$ for barred $k$, and ``(i leir)'' for ``(i bis).'' Its later ``(i)'' corresponds to a doubtful French (ii)/(iii), and the proof lacks the French point $y$. English 50 repeats the properness and surjectivity clauses, omits the French qualification about finding an open subset, and says ``to T and $\xi$'' for \emph{en $\xi$}. Its $H1$, four $g-1$ groups, and two $Y1$ groups remain on the baseline, alongside genuine $g_1$ indices. The query after $h$ corresponds to canceled fragments on French 44; English $\psi$ differs from the eventually secure French $\mu$. English 50 bars the second $Y1$ in the closure sentence, against the first $Y^1$ in French 45. Sources: \href{https://grothendieck.umontpellier.fr/web/viewer.html?file=../24.pdf\#page=44}{cote 24, sites 43--50}; \href{https://grothendieck.umontpellier.fr/web/viewer.html?file=../25.pdf\#page=38}{cote 25, sites 37--45}.

\paragraph{Blass's Section 8: singularities and linear subsystems (sites 51--62).}
The continuation completes 8.9, discusses the singular and supersingular loci and 8.10--8.18, then returns to linear subsystems and the projected construction. English 51 prints baseline $U1$ for French $U'$, omits the explicit isomorphism induced by $g_1$, and has equality in the dimension bound where French 45 retains a doubtful replacement by $\leq$. Its query about the singular locus points to the French definition by the vanishing of a section $D$ of an invertible Module. English 52's queries refer to French 46's tangent at $x$, its revised singular-locus notation, and the repeated geometric-singularity phrase; the English adds ``double'' to ``ordinary points,'' and its clear degree 4 leaves the French 4/5 doubt unresolved. English 53 retains its literal $W$ formula and index levels against the French handwritten omega. Its three separate gaps point to French 47's top note and the inserted dimension-one example; the French marginal definition of \emph{réglé} is absent. English 54's hyperplane query corresponds to French 47's reorganization note, and it omits the live premise that $X$ satisfies geometric $(R_k)$. In 8.13 it prints capital $K$ and the bound $\leq K-1$, against French 48's $\leq k+1$, while leaving out the marginal geometric-singular-locus and dimension notes. English 55 retains a repeated depth expression, $\operatorname{prof}xX\geq2\operatorname{prof}_xX$, against the single French bound, and omits ``geometrically'' before ``normal'' in 8.15. English 56 leaves the codimension bound blank where French 49 has $\geq2$, retains a narrow left angle after $Z$, and prints $K+1$ against a doubtful French letter or digit. English 57's four gaps correspond individually to French 50--51; 8.17(a) omits the French qualification beginning \emph{Il est possible que} and the geometric-irreducibility phrase, prints $f(x)$, and calls the tangent hyperplanes ``hyperplane sections.'' English 58 preserves the French-language query about \emph{réglé}, ``non???'' before ``singular,'' and a further gap of uncertain attachment. It omits $d$ before the dimension bound, uses a slash for restriction, and raises 1 where French 52 has a prime. English 59 omits two image-closure bars and dual marks, prints capital $X\in T$, and says ``and not only'' where the French says \emph{mais seulement}; its translation alternatives and ``Fr p50'' stay literal. English 60 preserves further alternatives and French glosses, omits several dual, zero, and $\xi$ marks, prints $f(x)\ C\ L0$ against French 53's $f(X)\not\subset L^o$, and uses capital $U$ and raised 1 in the diagrams. English 61 prints $\dim f(X)=2$ where French 53 has a not-equal sign and a doubtful 2/1, omits field bars and the zero on $L$, starts the field expression with capital $K$, says ``dimension at most one'' for \emph{dimension une de plus}, and leaves the target image unbarred. The final English 62 remark retains lowercase $x$ where French 54 has $X$, and its ``Bertinisque'' spelling. Every remaining query is kept with its recorded French location. Sources: \href{https://grothendieck.umontpellier.fr/web/viewer.html?file=../24.pdf\#page=52}{cote 24, sites 51--62}; \href{https://grothendieck.umontpellier.fr/web/viewer.html?file=../25.pdf\#page=46}{cote 25, sites 45--54}.

\paragraph{Blass's Section 9 (sites 63--66).}
The translation follows the Veronese-type constructions and the degree-dependent results through the proof of 9.3, then preserves the unfinished 9.4 and 9.5. English 63 retains the missing relation in the opening direct-image/Sym expression, also absent in French 55; its ordinary and lowered stars and its question marks are separately preserved. Its literal $O\ P(n)$ uses ordinary baseline letters, and the following $P(n)$ lacks the French dual mark. The wide gap after ``embarrassing f'', the French phrase without accents, and ``or (admits)'' correspond to French 55's \emph{de f à certains égards} and \emph{ou n'admet qu'un}. English 64 repeats an N.B. clause that appears once in French 55--56. In 9.3(c), ``where'' corresponds to the inserted French \emph{ou}; the field is unbarred, the first $Y(n)$ lacks the attached index, and the last exponent has a closing parenthesis without its opening mate. English 65's queries correspond to French 57's clause about $t_s$, \emph{on fait pareil}, and a live/canceled-text boundary. The affine-model sentence uses $f(X)$ and ends its coordinates at $T_s$, against French 56's $f(x)$ and a doubtful $T|X$ with index $r$. ``Move margin'', capital $K$, and the capital C in the image expression are literal. The $i/t_s$ mismatch is shared with French 57. English 66 prints a clear $Q/R$ where the French retains $Q/P$ doubt. Its blank Corollary 9.4 and all the 9.5 questions point to French 57's unfinished headings, conditions, and notes. Neither witness supplies a completed statement there. Sources: \href{https://grothendieck.umontpellier.fr/web/viewer.html?file=../24.pdf\#page=64}{cote 24, sites 63--66}; \href{https://grothendieck.umontpellier.fr/web/viewer.html?file=../25.pdf\#page=56}{cote 25, sites 55--57}.

\paragraph{Blass's Section 10 (sites 66--71).}
The section translates the discussion of sections by linear subvarieties, the regular-immersion construction, and the planning comments around 10.2--10.4. English 67 keeps ``not meeting k'' and adds ``(Z AG's error P.B.)'' with a handwritten question; French 58 indeed has $k$, so this remains Blass's proposed correction. Several closure bars and dual marks are absent, the English uses raised 1 for a French prime, and its clear $f(T)$ leaves the French $F|T$ unresolved. English 68 keeps ``pull out of find out'' and two N.B. closures. Three sheaf occurrences have the narrow digit $0$ where other occurrences on that page use round capital $O$; the underlines, product letter x, and typed dual v's remain as printed. English 69 refers to ``No.1'' where French 60 inserts 12, repeats overlapping opening clauses, and retains the mystery comment, the gap before 5.3, and the handwritten request. The corresponding French reference remains doubtful between S.3 and 5.3. Its ``only'' replaces \emph{surtout}, ``which according to 10.2'' replaces \emph{qu'après 10.2}, and the codimension-two wording lacks the uncertain sign visible in French 60. English 70 preserves the page reference ``[58]'', the extra qualification ``to some extent'', and ``except'' in a sentence whose French has repeated \emph{sont}. The query about a regular-immersion section points to French 61's \emph{un n° en rabiot}. The typescript then says ``Proposition10.3 is crossed out. AskAG if that is his intention.'' It omits the proposition body, the exercise paragraph, and the dimension-at-most-one transition present in checked French 61; the claim about cancellation is attributed to Blass. English 71 asks for the translation of ``sorites'' and adds ``(facts)'', at French 62's doubtful \emph{sorties|sorttes}; its handwritten translation question remains partly uncertain. Other marked alternatives concern taking the construction in practice, revisiting it, $X$ or $S$, and closed immersions. ``Forefathers'' renders \emph{Papa}, the phrase \emph{au dessus de s} is omitted, and $n-1$ corresponds to a doubtful French $t|r-1$. Sources: \href{https://grothendieck.umontpellier.fr/web/viewer.html?file=../24.pdf\#page=67}{cote 24, sites 66--71}; \href{https://grothendieck.umontpellier.fr/web/viewer.html?file=../25.pdf\#page=59}{cote 25, sites 58--62}.

\paragraph{Blass's Section 11 (sites 71--74).}
The final translated section develops Grassmannians and linear subspaces of projective bundles, corresponding to French 62--65. English 71 prints $nE$, a plain rank-one $G$, and $P(nE)$ without the French exterior-power signs, and leaves several sheaf letters ununderlined. English 72 places handwritten $r$ above typed $(n)$ with only $n$ inside the parentheses; it repeats two representability lines where French 63 has one clause. It twice prints capital $S$ in the Grassmannian index, says ``recovered from'' for \emph{recouvert par}, and omits ``vector'' from ``fiber bundle''. The sequence endpoints are round capital O rather than the French numerical zeros; all four English connectors have arrowheads, where the two middle French strokes do not. The Hom subscript is a narrow digit 0 followed by lowercase s on one lowered tier, against French underlined O with lowered S. English 73 omits the characterization by closed subpreschemes that are smooth over S and everywhere of the stated relative dimension, joining directly to projective degree one. Its ``empty is n=0'' corresponds to \emph{et vide si n=0}. The first English dual v is clear where the corresponding first French dual mark is absent; the other typed v's and raised 1's remain at their observed levels. English 74's many queries point individually to French 65: ``again'', ``by which of course we mean'', the G/Q gap, the sorites gloss, ``(only when)until'', the \emph{commode} gloss, and the two illegibility requests all remain. The points-of-$L$ phrase lacks the French raised o; the final $E$ index has a baseline 1 after lowered S. The capital C and ``in so far as'' replace a French hooked arrow and \emph{ssi}, respectively. The handwritten star/bullet indices on Grass and the coproduct sign before the indexed Grassmannians are absent in the English. Its unmatched N.B. bracket and final handwritten request to A.G. or Deligne are preserved. Sources: \href{https://grothendieck.umontpellier.fr/web/viewer.html?file=../24.pdf\#page=72}{cote 24, sites 71--74}; \href{https://grothendieck.umontpellier.fr/web/viewer.html?file=../25.pdf\#page=63}{cote 25, sites 62--65}.

\paragraph{The two final sheets (sites 75--76).}
The short handwritten French note on site 75 refers to having several important activities at once, one mathematical and another nonmathematical but worthwhile. Its initial \emph{Idem} and \emph{aussi} remain doubtful. The interlinear \emph{importantes} is connected to the gap after \emph{activités} by a large pen leader, whose right arm underlines the inserted word. Mirrored typing visible through the sheet belongs to the reverse and is described in a note rather than transcribed a second time. Site 76 is that separate typed CV ending, signed A. Grothendieck, with a statement about \emph{Survivre}, survival, and scientific activity, followed by a list of principal publications. Its two marginal pen brackets, separate short stroke, and folio number remain visible in the diplomatic transcription. The dark wedge below \emph{parmi} is a hole in the paper, established by the corresponding mirrored feature on site 75. This sheet's inclusion in the diplomatic volume preserves the archive sequence while its stray-sheet label explains its position. Sources: \href{https://grothendieck.umontpellier.fr/web/viewer.html?file=../24.pdf\#page=76}{cote 24, site 75}; \href{https://grothendieck.umontpellier.fr/web/viewer.html?file=../24.pdf\#page=77}{cote 24, site 76}.
\par\typeout{EGA-SOURCE-END guide-blass PDF-PAGE \thepage}

\par\typeout{EGA-SOURCE-BEGIN guide-published-references PDF-PAGE \thepage}
\pdfbookmark[1]{Published EGA references}{source:guide-published-references}
\section*{Published references and the manuscript's numbering}
The numbered references in the commentary preserve those read in the prenotes. The published EGA IV provides the following verified points of comparison: the codimension criterion for $(S_k)$ at 5.7.4; the fiber-property lists at 12.1.1 and 12.1.6; the proper-family discussion at 12.2.1; principal parts at 16.3.1--16.3.2; and regular and quasi-regular immersions at 16.9.1--16.9.3. The manuscript's use of 12.1.4 for a list requires the note given under IV 20.5: published 12.1.4 concerns fiber dimension. The manuscript citations 15.2.2, 15.6.3 and IV 23.9.2 are reported as written; no replacement number is supplied without an established concordance. Published EGA IV was consulted for this bibliographical context. Every reading of the prenotes and of Blass's typescript was decided from the supplied scans.

\begin{itemize}
\item A. Grothendieck, with the collaboration of J. Dieudonn\'e, \emph{\'El\'ements de g\'eom\'etrie alg\'ebrique, IV. \'Etude locale des sch\'emas et des morphismes de sch\'emas, Seconde partie}, Publications Math\'ematiques de l'IH\'ES \textbf{24} (1965), 5--231. \href{https://www.numdam.org/articles/10.1007/BF02684322/}{doi:10.1007/BF02684322}.
\item A. Grothendieck, with the collaboration of J. Dieudonn\'e, \emph{\'El\'ements de g\'eom\'etrie alg\'ebrique, IV. \'Etude locale des sch\'emas et des morphismes de sch\'emas, Troisi\`eme partie}, Publications Math\'ematiques de l'IH\'ES \textbf{28} (1966), 5--255. \href{https://www.numdam.org/articles/10.1007/BF02684343/}{doi:10.1007/BF02684343}.
\item A. Grothendieck, with the collaboration of J. Dieudonn\'e, \emph{\'El\'ements de g\'eom\'etrie alg\'ebrique, IV. \'Etude locale des sch\'emas et des morphismes de sch\'emas, Quatri\`eme partie}, Publications Math\'ematiques de l'IH\'ES \textbf{32} (1967), 5--361. \href{https://www.numdam.org/articles/10.1007/BF02732123/}{doi:10.1007/BF02732123}.
\end{itemize}
\par\typeout{EGA-SOURCE-END guide-published-references PDF-PAGE \thepage}

\end{document}
